

\setcounter{chapter}{3}

\chapter{群作用}

在本章中我们考察群作用在一个集合上的一些推论。当一个对象作用于另一个对象时，可以从两者获得大量信息，这是数学中一个重要而反复出现的观念。随着群所作用的集合被赋予更多结构（例如群作用于群，或群作用于向量空间（将在第 18 章讨论）），关于群结构的更多信息就变得可用。对群作用的这一研究在本章达到高潮，集中体现为西罗定理的证明以及由此产生的例子和分类。

作用的概念将在我们研究模、向量空间、矩阵的标准形和伽罗瓦理论时反复出现，是本书中基本的统一主题之一。

\section{群作用与置换表示}

本节我们给出群作用的基本理论，然后把这套理论应用于作用在 \(\{1, 2, \ldots, n\}\) 上的 \(S_n\) 的子群，以证明 \(S_n\) 的每个元素都有唯一的循环分解。在第 2、3 节中，我们把一般理论应用于另外两个具体的群作用，以导出一些重要结果。

设群 \(G\) 作用在非空集合 \(A\) 上。回忆第 1.7 节，对每个 \(g \in G\)，由
\[
\sigma_g(a) = g \cdot a
\]
定义的映射 \(\sigma_g : A \to A\) 是 \(A\) 的一个置换。我们在第 1.7 节还看到，对 \(G\) 在 \(A\) 上的作用有与之相伴的同态
\[
\varphi : G \to S_A, \quad \text{由 } \varphi(g) = \sigma_g \text{ 定义},
\]
称为与该作用相伴的\textbf{置换表示}。回忆第 1.7 节和第 2.2 节引入的关于群作用的一些附加术语。

\begin{definition}
（1）作用的\textbf{核}是在 \(A\) 的每个元素上都平凡地作用的 \(G\) 中元素的集合：\(\{g \in G \mid g \cdot a = a \text{ 对所有 } a \in A\}\).

（2）对每个 \(a \in A\)，\(a\) 在 \(G\) 中的\textbf{稳定子}是固定元素 \(a\) 的 \(G\) 中元素的集合：\(\{g \in G \mid g \cdot a = a\}\)，并记作 \(G_a\).

（3）若作用的核是单位元，则称该作用是\textbf{忠实的}。
\end{definition}

注意，作用的核恰好与相伴置换表示的核相同；特别地，核是 \(G\) 的正规子群。两个群元素在 \(A\) 上诱导相同的置换，当且仅当它们位于核的同一陪集中（当且仅当它们位于置换表示的同一纤维中）。特别地，\(G\) 在 \(A\) 上的作用也可以看作商群 \(G/\ker \varphi\) 在 \(A\) 上的忠实作用。回忆第 2.2 节，\(A\) 中元素 \(a\) 在 \(G\) 中的稳定子是 \(G\) 的子群。若 \(a\) 是 \(A\) 的固定元素，则作用的核包含在稳定子 \(G_a\) 中，因为作用的核是稳定每个点的 \(G\) 中元素的集合，即 \(\bigcap_{a \in A} G_a\).

\begin{example}
（1）设 \(n\) 是正整数。群 \(G = S_n\) 通过 \(\sigma \cdot i = \sigma(i)\)（对所有 \(i \in \{1, \ldots, n\}\)）作用在集合 \(A = \{1, 2, \ldots, n\}\) 上。与该作用相伴的置换表示是恒等映射 \(\varphi : S_n \to S_n\). 这个作用是忠实的，且对每个 \(i \in \{1, \ldots, n\}\)，稳定子 \(G_i\)（所有固定 \(i\) 的置换构成的子群）同构于 \(S_{n-1}\)（参见第 3.2 节习题 15）。

（2）设 \(G = D_8\) 作用在由正方形的四个顶点组成的集合 \(A\) 上。按第 1.2 节图 2 的方式，把这些顶点顺时针标为 1, 2, 3, 4. 设 \(r\) 是把正方形顺时针旋转 \(\pi/2\) 弧度的旋转，\(s\) 是关于经过顶点 1 和 3 的直线的反射。于是由 \(r\) 和 \(s\) 给出的顶点置换是
\[
\sigma_r = (1\,2\,3\,4) \quad \text{和} \quad \sigma_s = (2\,4).
\]
注意，由于置换表示是同态，与 \(sr\) 对应的四个顶点的置换是 \(\sigma_{sr} = \sigma_s\sigma_r = (1\,4)(2\,3)\). \(D_8\) 在正方形四个顶点上的作用是忠实的，因为只有恒等对称固定所有四个顶点。任意顶点的稳定子是 \(D_8\) 中由关于经过该顶点和正方形中心的直线的反射生成的 2 阶子群（例如，顶点 1 的稳定子是 \(\langle s \rangle\)）。

（3）像前一个例子那样给正方形的四个顶点标号，现在设 \(A\) 是由不相邻（相对）顶点组成的无序对的集合：\(A = \{\{1, 3\}, \{2, 4\}\}\). 于是 \(D_8\) 也作用在这个集合 \(A\) 上，因为正方形的每个对称把一对相对顶点送到一对相对顶点。旋转 \(r\) 交换对 \(\{1, 3\}\) 和 \(\{2, 4\}\)；反射固定这两对相对顶点。于是若把这两对 \(\{1, 3\}\) 和 \(\{2, 4\}\) 分别标为 1 和 2，则由 \(r\) 和 \(s\) 给出的 \(A\) 的置换是
\[
\sigma_r = (1\,2) \quad \text{和} \quad \sigma_s = \text{恒等置换}.
\]
这个 \(D_8\) 的作用不忠实：它的核是 \(\langle s, r^2 \rangle\). 此外，对每个 \(\alpha \in A\)，\(\alpha\) 在 \(D_8\) 中的稳定子都等于作用的核。

（4）像例 2 那样给正方形的四个顶点标号，现在设 \(A\) 是如下无序顶点对的集合：\(\{\{1, 2\}, \{3, 4\}\}\). 群 \(D_8\) 不作用在这个集合 \(A\) 上，因为 \(\{1, 2\} \in A\) 但 \(r \cdot \{1, 2\} = \{2, 3\} \notin A\).
\end{example}

作用与映入对称群的同态之间的关系可以反过来。也就是说，给定任意非空集合 \(A\) 和群 \(G\) 到 \(S_A\) 的任意同态 \(\varphi\)，我们通过定义
\[
g \cdot a = \varphi(g)(a)
\]
（对所有 \(g \in G\) 和所有 \(a \in A\)）得到 \(G\) 在 \(A\) 上的一个作用。这个作用的核与 \(\ker \varphi\) 相同。与该作用相伴的置换表示恰好就是给定的同态 \(\varphi\). 这证明了下面的结果。

\begin{proposition}\label{pro:4.1}
对任意群 \(G\) 和任意非空集合 \(A\)，\(G\) 在 \(A\) 上的作用与 \(G\) 到 \(S_A\) 的同态之间存在一个双射。
\end{proposition}

鉴于命题 \ref{pro:4.1}，置换表示的定义可以重新表述。

\begin{definition}
若 \(G\) 是群，则 \(G\) 的\textbf{置换表示}是 \(G\) 到某个非空集合 \(A\) 的对称群 \(S_A\) 的任意同态。我们说 \(G\) 在 \(A\) 上的给定作用\textbf{提供}或\textbf{诱导}与之相伴的 \(G\) 的置换表示。
\end{definition}

我们可以把置换表示看作线性变换的矩阵表示的类比。在 \(A\) 是含 \(n\) 个元素的有限集合的情形，有 \(S_A \cong S_n\)（参见第 1.6 节），所以通过固定 \(A\) 中元素的一个标号，我们可以把置换看作群 \(S_n\) 的元素（这正是我们在上面例 2 和例 3 中所做的），就像固定向量空间的一个基使我们能把线性变换看作矩阵一样。

现在我们证明一个关于群作用的组合结果，当我们在后续各节把它应用于具体的群作用时，它会有重要的推论。

\begin{proposition}\label{pro:4.2}
设群 \(G\) 作用在非空集合 \(A\) 上。由
\[
a \sim b \quad \text{当且仅当} \quad a = g \cdot b \text{ 对某个 } g \in G
\]
定义的 \(A\) 上的关系是等价关系。对每个 \(a \in A\)，包含 \(a\) 的等价类中元素的个数为 \(|G : G_a|\)，即 \(a\) 的稳定子的指数。
\end{proposition}

\begin{proof}
我们首先证明 \(\sim\) 是等价关系。由作用的公理 2，对所有 \(a \in A\) 有 \(1 \cdot a = a\)，即 \(a \sim a\)，关系是自反的。若 \(a \sim b\)，则对某个 \(g \in G\) 有 \(a = g \cdot b\)，于是
\[
g^{-1} \cdot a = g^{-1} \cdot (g \cdot b) = (g^{-1}g) \cdot b = 1 \cdot b = b,
\]
即 \(b \sim a\)，关系是对称的。最后，若 \(a \sim b\) 且 \(b \sim c\)，则对某个 \(g, h \in G\) 有 \(a = g \cdot b\)、\(b = h \cdot c\)，故
\[
a = g \cdot b = g \cdot (h \cdot c) = (gh) \cdot c,
\]
于是 \(a \sim c\)，关系是传递的。

为证明命题的最后一个论断，我们给出 \(G_a\) 在 \(G\) 中的左陪集与 \(a\) 的等价类中元素之间的一个双射。设 \(C_a\) 是 \(a\) 所在的类，则
\[
C_a = \{g \cdot a \mid g \in G\}.
\]
设 \(b = g \cdot a \in C_a\). 则 \(gG_a\) 是 \(G_a\) 在 \(G\) 中的一个左陪集。映射
\[
b = g \cdot a \longmapsto gG_a
\]
是从 \(C_a\) 到 \(G_a\) 在 \(G\) 中左陪集集合的映射。这个映射是满射，因为对任意 \(g \in G\)，元素 \(g \cdot a\) 都是 \(C_a\) 的元素。由于 \(g \cdot a = h \cdot a\) 当且仅当 \(h^{-1}g \in G_a\) 当且仅当 \(gG_a = hG_a\)，这个映射也是单射，因此是一个双射。这就完成了证明。
\end{proof}

由命题 \ref{pro:4.2}，群 \(G\) 作用在集合 \(A\) 上把 \(A\) 划分成在 \(G\) 作用下的不相交的等价类。这些类被赋予一个名称：

\begin{definition}
设群 \(G\) 作用在非空集合 \(A\) 上。

（1）等价类 \(\{g \cdot a \mid g \in G\}\) 称为包含 \(a\) 的\textbf{轨道}。

（2）若只有一个轨道，即若对任意两个元素 \(a, b \in A\) 都存在 \(g \in G\) 使 \(g \cdot a = b\)，则称 \(G\) 在 \(A\) 上的作用是\textbf{传递的}。
\end{definition}

\begin{example}
设群 \(G\) 作用在集合 \(A\) 上。

（1）若 \(G\) 平凡地作用在 \(A\) 上，则对所有 \(a \in A\) 有 \(G_a = G\)，且轨道就是 \(A\) 的元素。这个作用是传递的当且仅当 \(|A| = 1\).

（2）对称群 \(G = S_n\) 在它作为置换对 \(A = \{1, 2, \ldots, n\}\) 的通常作用下是传递的。注意任意点 \(i\) 在 \(G\) 中的稳定子在 \(S_n\) 中指数为 \(n = |A|\).

（3）当群 \(G\) 作用在集合 \(A\) 上时，\(G\) 的任意子群也作用在 \(A\) 上。若 \(G\) 在 \(A\) 上传递，则 \(G\) 的子群不必在 \(A\) 上传递。例如，若 \(G = \langle (1\,2), (3\,4) \rangle \leq S_4\)，则 \(G\) 在 \(\{1, 2, 3, 4\}\) 上的轨道是 \(\{1, 2\}\) 和 \(\{3, 4\}\)，且 \(G\) 中没有元素把 2 送到 3. 下面关于循环分解的讨论表明，当 \(\langle \sigma \rangle\) 是 \(S_n\) 的任意循环子群时，\(\langle \sigma \rangle\) 的轨道由 \(\sigma\) 的循环分解中各个循环里出现的数集组成（例如，\(\langle (1\,2)(3\,4\,5) \rangle\) 的轨道是 \(\{1, 2\}\) 和 \(\{3, 4, 5\}\)）。

（4）群 \(D_8\) 传递地作用在正方形的四个顶点上，任意顶点的稳定子是 2 阶（指数为 4）的子群，由关于经过该点的对称轴的反射生成。

（5）群 \(D_8\) 也传递地作用在两对相对顶点组成的集合上。在这个作用中，任意点的稳定子是 \(\langle s, r^2 \rangle\)（指数为 2）。
\end{example}

\subsection*{循环分解}

我们现在证明对称群 \(S_n\) 的每个元素都有第 1.3 节所述的唯一循环分解。设 \(A = \{1, 2, \ldots, n\}\)，设 \(\sigma\) 是 \(S_n\) 的元素，令 \(G = \langle \sigma \rangle\). 则 \(\langle \sigma \rangle\) 作用在 \(A\) 上，于是由命题 \ref{pro:4.2}，它把 \(\{1, 2, \ldots, n\}\) 划分成唯一一组（不相交的）轨道。设 \(O\) 是其中一个轨道，\(x \in O\). 对 \(A = O\) 应用命题 \ref{pro:4.2}（的证明），我们看到 \(G_x\) 在 \(G\) 中的左陪集与 \(O\) 的元素之间存在双射，由
\[
g \cdot x \longmapsto gG_x
\]
明确给出。由于 \(G\) 是循环群，\(G_x \trianglelefteq G\)，且 \(G/G_x\) 是 \(d\) 阶循环群，其中 \(d\) 是使 \(\sigma^d \in G_x\) 的最小正整数（参见第 3.1 节命题 7 后的例 2）。又 \(d = |G : G_x| = |O|\). 于是 \(G_x\) 在 \(G\) 中的不同陪集是
\[
G_x, \sigma G_x, \sigma^2 G_x, \ldots, \sigma^{d-1}G_x.
\]
这表明 \(O\) 的不同元素是
\[
x, \sigma(x), \sigma^2(x), \ldots, \sigma^{d-1}(x).
\]
按这种方式给 \(O\) 的元素排序表明 \(\sigma\) 循环地置换 \(O\) 的元素，也就是说，在大小为 \(d\) 的轨道上，\(\sigma\) 起一个 \(d\)-循环的作用。这证明了对每个 \(\sigma \in S_n\) 都存在循环分解。

\(\langle \sigma \rangle\) 的轨道由 \(\sigma\) 唯一确定。唯一的自由度在于列出轨道的顺序。在每个轨道 \(O\) 内，我们可以任取一个元素作为代表元开始。取 \(\sigma^i(x)\) 而不是 \(x\) 作为初始代表元，只是把 \(O\) 的元素按顺序
\[
\sigma^i(x), \sigma^{i+1}(x), \ldots, \sigma^{d-1}(x), x, \sigma(x), \ldots, \sigma^{i-1}(x)
\]
列出，即原来列表的一个循环置换（向前移动 \(i-1\) 项）。由此可知，上面的循环分解在循环的重新排列以及每个循环内整数的循环置换的意义下是唯一的。

对称群的子群称为\textbf{置换群}。对 \(S_n\) 的任意子群 \(G\)，\(G\) 的轨道是指它在 \(\{1, 2, \ldots, n\}\) 上的轨道。\(S_n\) 中元素 \(\sigma\) 的轨道是指群 \(\langle \sigma \rangle\) 的轨道（即构成其循环分解中各个循环的整数集合）。

下面的习题进一步说明如何从置换表示获得群论信息。

\subsection*{习题}
设 \(G\) 是群，\(A\) 是非空集合。

\begin{enumerate}
\item 设 \(G\) 作用在集合 \(A\) 上。证明：若 \(a, b \in A\) 且对某个 \(g \in G\) 有 \(b = g \cdot a\)，则 \(G_b = gG_a g^{-1}\)（\(G_a\) 是 \(a\) 的稳定子）。由此推出：若 \(G\) 传递地作用在 \(A\) 上，则作用的核是 \(\bigcap_{g \in G} gG_a g^{-1}\).

\item 设 \(G\) 是集合 \(A\) 上的置换群（即 \(G \leq S_A\)），\(\sigma \in G\)，\(a \in A\). 证明 \(\sigma G_a \sigma^{-1} = G_{\sigma(a)}\). 由此推出：若 \(G\) 传递地作用在 \(A\) 上，则 \(\bigcap_{\sigma \in G} \sigma G_a \sigma^{-1} = 1\).

\item 假设 \(G\) 是 \(S_A\) 的阿贝尔传递子群。证明对所有 \(\sigma \in G - \{1\}\) 和所有 \(a \in A\) 有 \(\sigma(a) \neq a\). 由此推出 \(|G| = |A|\).〔使用前一习题。〕

\item 设 \(S_3\) 作用在有序对的集合 \(Q = \{(i, j) \mid 1 \leq i, j \leq 3\}\) 上，作用为 \(\sigma((i, j)) = (\sigma(i), \sigma(j))\). 求 \(S_3\) 在 \(Q\) 上的轨道。对每个 \(\sigma \in S_3\) 求它在这个作用下的循环分解（即把 \(\sigma\) 看作 \(S_9\) 的元素时的循环分解——先给这九个有序对固定一个标号）。对 \(S_3\) 作用在这九个点上得到的每个轨道 \(O\)，取某个 \(a \in O\)，求 \(a\) 在 \(S_3\) 中的稳定子。

\item 对 (a) 和 (b) 分别重复前一习题，但把 \(S_3\) 作用在指定的集合上：

（a）27 个三元组 \(\{(i, j, k) \mid 1 \leq i, j, k \leq 3\}\) 的集合；

（b）\(\{1, 2, 3\}\) 的全部 7 个非空子集组成的集合 \(\mathcal{P}(\{1, 2, 3\}) - \{\varnothing\}\).

\item 如第 2.2 节习题 12 那样，设 \(R\) 是独立变量 \(x_1, x_2, x_3, x_4\) 的整系数多项式集合，并让 \(S_4\) 通过置换四个变量的下标作用于 \(R\)：
\[
\sigma \cdot p(x_1, x_2, x_3, x_4) = p(x_{\sigma(1)}, x_{\sigma(2)}, x_{\sigma(3)}, x_{\sigma(4)})
\]
对所有 \(\sigma \in S_4\).

（a）求 \(S_4\) 在 \(R\) 上的轨道中含 \(x_1 + x_2\) 的多项式（回忆第 2.2 节习题 12，这个多项式的稳定子阶为 4）。

（b）求 \(S_4\) 在 \(R\) 上的轨道中含 \(x_1x_2 + x_3x_4\) 的多项式（回忆第 2.2 节习题 12，这个多项式的稳定子阶为 8）。

（c）求 \(S_4\) 在 \(R\) 上的轨道中含 \((x_1 + x_2)(x_3 + x_4)\) 的多项式。

\item 设 \(G\) 是有限集合 \(A\) 上的传递置换群。\textbf{块}是 \(A\) 的非空子集 \(B\)，使得对所有 \(\sigma \in G\)，或者 \(\sigma(B) = B\)，或者 \(\sigma(B) \cap B = \varnothing\)（这里 \(\sigma(B)\) 是集合 \(\{\sigma(b) \mid b \in B\}\)）。

（a）证明：若 \(B\) 是包含 \(A\) 中元素 \(a\) 的块，则由
\[
G_B = \{\sigma \in G \mid \sigma(B) = B\}
\]
定义的集合 \(G_B\) 是 \(G\) 的包含 \(G_a\) 的子群。

（b）证明：若 \(B\) 是块，且 \(\sigma_1(B), \sigma_2(B), \ldots, \sigma_n(B)\) 是 \(B\) 在 \(G\) 的元素作用下的全部不同像，则它们构成 \(A\) 的一个划分。

（c）若集合 \(A\) 上的（传递）群 \(G\) 仅有的块是平凡的（即大小为 1 的集合和 \(A\) 本身），则称 \(G\) 是\textbf{本原的}。证明 \(S_4\) 在 \(A = \{1, 2, 3, 4\}\) 上本原。证明 \(D_8\) 作为正方形四个顶点上的置换群不是本原的。

（d）证明传递群 \(G\) 在 \(A\) 上本原，当且仅当对每个 \(a \in A\)，包含 \(G_a\) 的 \(G\) 的子群只有 \(G_a\) 和 \(G\)（即 \(G_a\) 是 \(G\) 的极大子群，参见第 2.4 节习题 16）。〔使用（a）。〕

\item 集合 \(A\) 上的传递置换群 \(G\) 称为\textbf{双重传递的}，若对任意（从而对所有）\(a \in A\)，子群 \(G_a\) 传递地作用在 \(A - \{a\}\) 上。

（a）证明对所有 \(n \geq 2\)，\(S_n\) 在 \(\{1, 2, \ldots, n\}\) 上双重传递。

（b）证明双重传递群是本原的。由此推出 \(D_8\) 在正方形 4 个顶点上的作用不是双重传递的。

\item 假设 \(G\) 传递地作用在有限集合 \(A\) 上，设 \(H\) 是 \(G\) 的正规子群。设 \(O_1, O_2, \ldots, O_r\) 是 \(H\) 在 \(A\) 上的不同轨道。

（a）证明 \(G\) 置换集合 \(O_1, O_2, \ldots, O_r\)，即对每个 \(g \in G\) 和每个 \(i \in \{1, \ldots, r\}\)，存在 \(j\) 使得 \(gO_i = O_j\)，其中 \(gO = \{g \cdot a \mid a \in O\}\)（即在习题 7 的记号下，集合 \(O_1, \ldots, O_r\) 是块）。证明 \(G\) 在 \(\{O_1, \ldots, O_r\}\) 上传递。由此推出 \(H\) 在 \(A\) 上的所有轨道有相同的基数。

（b）证明：若 \(a \in O_1\)，则 \(|O_1| = |H : H \cap G_a|\)，并证明 \(r = |G : HG_a|\).〔画出描述子群 \(H\) 和 \(G_a\) 的第二同构定理的次格。注意 \(H \cap G_a = H_a\).〕

\item 设 \(H\) 和 \(K\) 是群 \(G\) 的子群。对每个 \(x \in G\)，定义 \(x\) 在 \(G\) 中的 \(HK\) \textbf{双陪集}为集合
\[
HxK = \{hxk \mid h \in H,\ k \in K\}.
\]
（a）证明 \(HxK\) 是左陪集 \(x_1K, \ldots, x_nK\) 的并，其中 \(\{x_1K, \ldots, x_nK\}\) 是 \(xK\) 在 \(H\) 通过左乘作用在 \(K\) 的左陪集集合上时的轨道。

（b）证明 \(HxK\) 是 \(H\) 的若干右陪集的并。

（c）证明对所有 \(x, y \in G\)，\(HxK\) 与 \(HyK\) 或者相同或者不相交。证明 \(HK\) 双陪集的集合划分 \(G\).

（d）证明 \(|HxK| = |K| \cdot |H : H \cap xKx^{-1}|\).

（e）证明 \(|HxK| = |H| \cdot |K : K \cap x^{-1}Hx|\).
\end{enumerate}

\section{群通过左乘作用在自身上——凯莱定理}

本节中 \(G\) 是任意群，我们首先考虑 \(G\) 通过左乘作用在自身上（即 \(A = G\)）：
\[
g \cdot a = ga \quad \text{对所有 } g \in G,\ a \in G,
\]
其中 \(ga\) 表示两个群元素 \(g\) 与 \(a\) 在 \(G\) 中的乘积（若 \(G\) 用加法书写，则该作用写成 \(g \cdot a = g + a\)，称为左平移）。我们在第 1.7 节看到它满足群作用的两条公理。

当 \(G\) 是 \(n\) 阶有限群时，为描述这个作用所提供的置换表示，把 \(G\) 的元素用整数 \(1, 2, \ldots, n\) 标号是方便的。这样 \(G\) 的元素被列为 \(g_1, g_2, \ldots, g_n\)，且对每个 \(g \in G\)，置换 \(\sigma_g\) 可以描述为下标 \(1, 2, \ldots, n\) 的一个置换如下：
\[
\sigma_g(i) = j \quad \text{当且仅当} \quad g g_i = g_j.
\]
对群元素作不同的标号将给出 \(\sigma_g\) 作为 \(\{1, 2, \ldots, n\}\) 的置换的不同描述（参见习题）。

\begin{example}
设 \(G = \{1, a, b, c\}\) 是第 2.5 节写出乘法表的克莱因四元群。把群元素 \(1, a, b, c\) 分别用整数 \(1, 2, 3, 4\) 标号。在这个标号下，我们计算由群元素 \(a\) 左乘的作用诱导的置换 \(\sigma_a\)：
\[
a \cdot 1 = a1 = a \text{，故 } \sigma_a(1) = 2;
\]
\[
a \cdot a = aa = 1 \text{，故 } \sigma_a(2) = 1;
\]
\[
a \cdot b = ab = c \text{，故 } \sigma_a(3) = 4;
\]
\[
a \cdot c = ac = b \text{，故 } \sigma_a(4) = 3.
\]
在这个 \(G\) 的元素标号下，我们看到 \(\sigma_a = (1\,2)(3\,4)\). 在克莱因四元群通过左乘作用在自身上所相伴的置换表示中，类似地计算得
\[
a \mapsto \sigma_a = (1\,2)(3\,4), \quad b \mapsto \sigma_b = (1\,3)(2\,4), \quad c \mapsto \sigma_c = (1\,4)(2\,3),
\]
这就具体给出了在这个标号下与该作用相伴的置换表示 \(G \to S_4\).
\end{example}

容易看出（我们不久将在更一般的框架中证明），群通过左乘作用在自身上总是传递且忠实的，并且任意点的稳定子都是单位子群（这些事实对上面的例子可以直接验证）。

我们现在考虑群通过左乘作用在其元素集合上的一个推广。设 \(H\) 是 \(G\) 的任意子群，\(A\) 是 \(H\) 在 \(G\) 中所有左陪集的集合。定义 \(G\) 在 \(A\) 上的作用为
\[
g \cdot aH = gaH \quad \text{对所有 } g \in G,\ aH \in A,
\]
其中 \(gaH\) 是以 \(ga\) 为代表的左陪集。容易验证它满足群作用的两条公理，即 \(G\) 确实通过左乘作用在 \(H\) 的左陪集集合上。当 \(H\) 是 \(G\) 的单位子群这一特殊情形下，陪集 \(aH\) 就是 \(\{a\}\)，若把元素 \(a\) 与集合 \(\{a\}\) 等同，则这个在单位子群的左陪集上通过左乘的作用与 \(G\) 通过左乘作用在自身上是同一个作用。

当 \(H\) 在 \(G\) 中指数为有限数 \(m\) 时，为描述这个作用所提供的置换表示，把 \(H\) 的左陪集用整数 \(1, 2, \ldots, m\) 标号是方便的。这样 \(H\) 在 \(G\) 中的不同左陪集被列为 \(a_1H, a_2H, \ldots, a_mH\)，且对每个 \(g \in G\)，置换 \(\sigma_g\) 可以描述为下标 \(1, 2, \ldots, m\) 的一个置换如下：
\[
\sigma_g(i) = j \quad \text{当且仅当} \quad g a_i H = a_j H.
\]
对左陪集作不同的标号将给出 \(\sigma_g\) 作为 \(\{1, 2, \ldots, m\}\) 的置换的不同描述（参见习题）。

\begin{example}
设 \(G = D_8\)，\(H = \langle s \rangle\). 把不同左陪集 \(1H, rH, r^2H, r^3H\) 分别用整数 \(1, 2, 3, 4\) 标号。在这个标号下，我们计算由群元素 \(s\) 左乘作用在 \(H\) 的左陪集上诱导的置换：
\[
s \cdot 1H = sH = 1H \text{，故 } \sigma_s(1) = 1;
\]
\[
s \cdot rH = srH = r^3H \text{，故 } \sigma_s(2) = 4;
\]
\[
s \cdot r^2H = sr^2H = r^2H \text{，故 } \sigma_s(3) = 3;
\]
\[
s \cdot r^3H = sr^3H = rH \text{，故 } \sigma_s(4) = 2.
\]
在这个 \(H\) 的左陪集标号下我们得到 \(\sigma_s = (2\,4)\). 在 \(D_8\) 通过左乘作用在 \(\langle s \rangle\) 的左陪集上所相伴的置换表示中，类似地计算得 \(\sigma_r = (1\,2\,3\,4)\). 注意置换表示是同态，所以一旦确定了它在 \(D_8\) 的生成元上的值，它在任意其他元素上的值就确定了（例如 \(\sigma_{sr^2} = \sigma_s\sigma_r^2\)）。
\end{example}

\begin{theorem}\label{thm:4.3}
设 \(G\) 是群，\(H\) 是 \(G\) 的子群，\(G\) 通过左乘作用在 \(H\) 在 \(G\) 中的左陪集集合 \(A\) 上。设 \(\pi_H\) 是这个作用所提供的相伴置换表示。则

\begin{enumerate}[label=(\arabic*)]
\item \(G\) 传递地作用在 \(A\) 上；
\item 点 \(1H \in A\) 在 \(G\) 中的稳定子是子群 \(H\)；
\item 作用的核（即 \(\pi_H\) 的核）是 \(\bigcap_{x \in G} xHx^{-1}\)，且 \(\ker \pi_H\) 是包含在 \(H\) 中的 \(G\) 的最大正规子群。
\end{enumerate}
\end{theorem}

\begin{proof}
为看出 \(G\) 传递地作用在 \(A\) 上，设 \(aH\) 和 \(bH\) 是 \(A\) 的任意两个元素，令 \(g = ba^{-1}\). 则 \(g \cdot aH = (ba^{-1})aH = bH\)，所以 \(A\) 的两个任意元素 \(aH\) 和 \(bH\) 位于同一轨道中，这证明了（1）。对（2），点 \(1H\) 的稳定子按定义是 \(\{g \in G \mid g \cdot 1H = 1H\}\)，即 \(\{g \in G \mid gH = H\} = H\).

由 \(\pi_H\) 的定义，
\[
\begin{aligned}
\ker \pi_H &= \{g \in G \mid g xH = xH \text{ 对所有 } x \in G\} \\
&= \{g \in G \mid (x^{-1}gx)H = H \text{ 对所有 } x \in G\} \\
&= \{g \in G \mid x^{-1}gx \in H \text{ 对所有 } x \in G\} \\
&= \{g \in G \mid g \in xHx^{-1} \text{ 对所有 } x \in G\} = \bigcap_{x \in G} xHx^{-1},
\end{aligned}
\]
这证明了（3）的第一个论断。（3）的第二个论断来自：首先注意 \(\ker \pi_H \trianglelefteq G\) 且 \(\ker \pi_H \leq H\). 若现在 \(N\) 是包含在 \(H\) 中的任意正规子群，则对所有 \(x \in G\) 有 \(N = xNx^{-1} \leq xHx^{-1}\)，故
\[
N \leq \bigcap_{x \in G} xHx^{-1} = \ker \pi_H.
\]
这表明 \(\ker \pi_H\) 是包含在 \(H\) 中的 \(G\) 的最大正规子群。
\end{proof}

\begin{corollary}[凯莱定理]\label{cor:4.4}
每个群都同构于某个对称群的子群。若 \(G\) 是 \(n\) 阶群，则 \(G\) 同构于 \(S_n\) 的子群。
\end{corollary}

\begin{proof}
令 \(H = 1\) 并应用前一个定理，得到 \(G\) 到 \(S_G\) 的同态（这里我们把单位子群的陪集与 \(G\) 的元素等同）。由于这个同态的核包含在 \(H = 1\) 中，\(G\) 同构于它在 \(S_G\) 中的像。
\end{proof}

注意 \(G\) 同构于某个对称群的一个子群，而不是同构于完整的对称群本身。例如，我们给出了克莱因四元群与 \(S_4\) 的子群 \(\langle (1\,2)(3\,4), (1\,3)(2\,4) \rangle\) 的一个同构。回忆对称群的子群称为置换群，所以凯莱定理表明每个群都同构于一个置换群。由左乘作用在 \(G\) 的元素（\(H = 1\) 的陪集）上所提供的置换表示称为 \(G\) 的\textbf{左正则表示}。

人们可能会想，我们只要研究对称群的子群就可以更有效地研究所有群（并通过研究对所有 \(n\) 的 \(S_n\) 的子群来研究所有有限群）。仅靠这种方法在计算和理论上都不实用，因为要研究 \(n\) 阶群，我们必须在一个大得多的群 \(S_n\) 中工作（例如参见习题 7）。

从历史上看，有限群最早不是像我们所发展的那样在公理化框架下被研究，而是作为 \(S_n\) 的子群被研究的。因此凯莱定理证明了历史上的群概念与现代（公理化）的群概念是等价的。现代方法的一个优点是：在研究一个给定群时，我们不必把该群限制为某个特定对称群的子群（所以在某种意义下我们的群是“无坐标的”）。

下一个结果推广了我们关于指数为 2 的子群正规的结论。

\begin{corollary}\label{cor:4.5}
若 \(G\) 是 \(n\) 阶有限群，\(p\) 是整除 \(|G|\) 的最小素数，则 \(G\) 的任何指数为 \(p\) 的子群都正规。
\end{corollary}

\textbf{注：}一般地，\(n\) 阶群不必有指数为 \(p\) 的子群（例如，\(A_4\) 没有指数为 2 的子群）。

\begin{proof}
假设 \(H \leq G\) 且 \(|G : H| = p\). 设 \(\pi_H\) 是由 \(G\) 在 \(H\) 的左陪集集合上左乘所提供的置换表示，令 \(K = \ker \pi_H\)，并设 \(|H : K| = k\). 则 \(|G : K| = |G : H||H : K| = pk\). 由于 \(H\) 有 \(p\) 个左陪集，由第一同构定理 \(G/K\) 同构于 \(S_p\) 的一个子群（即 \(G\) 在 \(\pi_H\) 下的像）。由拉格朗日定理，\(pk = |G/K|\) 整除 \(p!\). 于是 \(k \mid \frac{p!}{p} = (p-1)!\). 但 \((p-1)!\) 的所有素因子都小于 \(p\)，而由 \(p\) 的极小性，\(k\) 的每个素因子都大于或等于 \(p\). 这迫使 \(k = 1\)，故 \(H = K \trianglelefteq G\)，完成证明。
\end{proof}

\subsection*{习题}
设 \(G\) 是群，\(H\) 是 \(G\) 的子群。

\begin{enumerate}
\item 设 \(G = \{1, a, b, c\}\) 是第 2.5 节写出乘法表的克莱因四元群。

（a）把 \(1, a, b, c\) 分别用整数 \(1, 2, 4, 3\) 标号，并证明在 \(G\) 到 \(S_4\) 的左正则表示下，非单位元如下映：
\[
a \mapsto (1\,2)(3\,4), \quad b \mapsto (1\,4)(2\,3), \quad c \mapsto (1\,3)(2\,4).
\]
（b）把 \(1, a, b, c\) 分别重新标号为 \(1, 4, 2, 3\)，并计算 \(G\) 的每个元素在 \(G\) 到 \(S_4\) 的左正则表示下的像。证明在这个标号下 \(G\) 在 \(S_4\) 中的像与（a）中 \(G\) 的像是同一个子群（尽管在两种不同标号下各个非单位元映到不同的置换）。

\item 把 \(S_3\) 的元素列为 \(1, (1\,2), (2\,3), (1\,3), (1\,2\,3), (1\,3\,2)\) 并分别用整数 \(1, 2, 3, 4, 5, 6\) 标号。写出 \(S_3\) 的每个元素在 \(S_3\) 到 \(S_6\) 的左正则表示下的像。

\item 设 \(r\) 和 \(s\) 是 8 阶二面体群的通常生成元。

（a）把 \(D_8\) 的元素列为 \(1, r, r^2, r^3, s, sr, sr^2, sr^3\) 并分别用整数 \(1, 2, \ldots, 8\) 标号。写出 \(D_8\) 的每个元素在 \(D_8\) 到 \(S_8\) 的左正则表示下的像。

（b）把 \(D_8\) 的这一组元素分别重新用整数 \(1, 3, 5, 7, 2, 4, 6, 8\) 标号，并在这个新标号下重新计算 \(D_8\) 的每个元素在左正则表示下的像。证明（a）和（b）中得到的 \(S_8\) 的两个子群不同。

\item 用 \(Q_8\) 的左正则表示构造 \(S_8\) 的两个元素，使它们生成 \(S_8\) 的一个同构于四元数群 \(Q_8\) 的子群。

\item 设 \(r\) 和 \(s\) 是 8 阶二面体群的通常生成元，\(H = \langle s \rangle\). 把 \(H\) 在 \(D_8\) 中的左陪集列为 \(1H, rH, r^2H, r^3H\).

（a）分别用整数 \(1, 2, 3, 4\) 给这些陪集标号。写出 \(D_8\) 的每个元素在由 \(D_8\) 左乘作用在 \(D_8\) 中 \(H\) 的 4 个左陪集集合上所得表示 \(\pi_H\) 下（到 \(S_4\)）的像。由此推出这个表示是忠实的（即所得的 \(S_4\) 的元素构成一个同构于 \(D_8\) 的子群）。

（b）把陪集表分别重新用整数 \(1, 3, 2, 4\) 标号，重复（a）。证明由这个标号得到的置换构成的 \(S_4\) 的子群与（a）中得到的子群不同。

（c）设 \(K = \langle sr \rangle\)，把 \(K\) 在 \(D_8\) 中的陪集列为 \(1K, rK, r^2K, r^3K\)，并用整数 \(1, 2, 3, 4\) 给它们标号。证明在这个标号下，由左乘作用在 \(K\) 的陪集上所得表示下 \(D_8\) 的像与（a）中是 \(S_4\) 的同一个子群（尽管子群 \(H\) 与 \(K\) 不同，且 \(D_8\) 的某些元素在两个同态下映到不同的置换）。

\item 设 \(r\) 和 \(s\) 是 8 阶二面体群的通常生成元，\(N = \langle r^2 \rangle\). 把 \(N\) 在 \(D_8\) 中的左陪集列为 \(1N, rN, sN, srN\). 分别用整数 \(1, 2, 3, 4\) 给这些陪集标号。写出 \(D_8\) 的每个元素在由 \(D_8\) 左乘作用在 \(D_8\) 中 \(N\) 的 4 个左陪集集合上所得表示 \(\pi_N\) 下（到 \(S_4\)）的像。由此推出这个表示不忠实，并证明 \(\pi_N(D_8)\) 同构于克莱因四元群。

\item 设 \(Q_8\) 是 8 阶四元数群。

（a）证明 \(Q_8\) 同构于 \(S_8\) 的一个子群。

（b）证明 \(Q_8\) 不同构于任何 \(n \leq 7\) 的 \(S_n\) 的子群。〔若 \(Q_8\) 作用在任意 7 阶集合 \(A\) 上，证明任意点 \(a \in A\) 的稳定子必包含子群 \(\langle -1 \rangle\).〕

\item 证明：若 \(H\) 有有限指数 \(n\)，则存在 \(G\) 的正规子群 \(K\) 使得 \(K \leq H\) 且 \(|G : K| \leq n!\).

\item 证明：若 \(p\) 是素数，\(G\) 是 \(p^\alpha\) 阶群（\(\alpha \in \mathbb{Z}^+\)），则 \(G\) 的每个指数为 \(p\) 的子群都正规。由此推出每个 \(p^\alpha\) 阶群都有 \(p^{\alpha-1}\) 阶正规子群。

\item 证明每个 6 阶非阿贝尔群都有 2 阶非正规子群。用它分类 6 阶群。〔构造一个到 \(S_3\) 的单同态。〕

\item 设 \(G\) 是有限群，\(\pi : G \to S_G\) 是左正则表示。证明：若 \(x\) 是 \(G\) 的 \(n\) 阶元素且 \(|G| = mn\)，则 \(\pi(x)\) 是 \(m\) 个 \(n\)-循环的乘积。由此推出 \(\pi(x)\) 是奇置换当且仅当 \(|x|\) 是偶数且 \(\frac{|G|}{|x|}\) 是奇数。

\item 设 \(G\) 和 \(\pi\) 如前一个习题。证明：若 \(\pi(G)\) 含有一个奇置换，则 \(G\) 有指数为 2 的子群。〔使用第 3.3 节习题 3。〕

\item 证明：若 \(|G| = 2k\) 且 \(k\) 是奇数，则 \(G\) 有指数为 2 的子群。〔用柯西定理构造一个 2 阶元素，然后用前面两个习题。〕

\item 设 \(G\) 是合数阶 \(n\) 的有限群，且具有如下性质：对整除 \(n\) 的每个正整数 \(k\)，\(G\) 都有 \(k\) 阶子群。证明 \(G\) 不是单群。
\end{enumerate}

\section{群通过共轭作用在自身上——类方程}

本节中 \(G\) 是任意群，我们首先考虑 \(G\) 通过共轭作用在自身上（即 \(A = G\)）：
\[
g \cdot a = gag^{-1} \quad \text{对所有 } g \in G,\ a \in G,
\]
其中 \(gag^{-1}\) 照常在群 \(G\) 中计算。这个定义满足群作用的两条公理，因为
\[
g_1 \cdot (g_2 \cdot a) = g_1 \cdot (g_2 a g_2^{-1}) = g_1(g_2 a g_2^{-1})g_1^{-1} = (g_1 g_2)a(g_1 g_2)^{-1} = (g_1 g_2) \cdot a
\]
且
\[
1 \cdot a = 1a1^{-1} = a
\]
对所有 \(g_1, g_2 \in G\) 和所有 \(a \in G\).

\begin{definition}
若存在 \(g \in G\) 使 \(b = gag^{-1}\)（即当且仅当它们位于 \(G\) 通过共轭作用在自身上所得的同一轨道中），则称 \(G\) 的两个元素 \(a\) 和 \(b\) 在 \(G\) 中\textbf{共轭}。\(G\) 通过共轭作用在自身上所得的轨道称为 \(G\) 的\textbf{共轭类}。
\end{definition}

\begin{example}
（1）若 \(G\) 是阿贝尔群，则 \(G\) 通过共轭作用在自身上是平凡作用：对所有 \(g, a \in G\) 有 \(g \cdot a = a\)，且对每个 \(a \in G\)，\(a\) 的共轭类是 \(\{a\}\).

（2）若 \(|G| > 1\)，则与左乘作用不同，\(G\) 不通过共轭传递地作用在自身上，因为 \(\{1\}\) 总是一个共轭类（即这个作用的一个轨道）。更一般地，单元素子集 \(\{a\}\) 是共轭类当且仅当对所有 \(g \in G\) 有 \(gag^{-1} = a\)，当且仅当 \(a\) 在 \(G\) 的中心中。

（3）在 \(S_3\) 中可以直接算出共轭类是 \(\{1\}\)、\(\{(1\,2), (1\,3), (2\,3)\}\) 和 \(\{(1\,2\,3), (1\,3\,2)\}\). 我们不久将发展出更容易计算共轭类的技巧，特别是在对称群中。
\end{example}

正如群通过左乘作用在自身上可以推广一样，共轭作用也可以推广。若 \(S\) 是 \(G\) 的任意子集，定义
\[
gSg^{-1} = \{gsg^{-1} \mid s \in S\}.
\]
群 \(G\) 通过定义 \(g \cdot S = gSg^{-1}\)（对任意 \(g \in G\)、\(S \in \mathcal{P}(G)\)）作用在自身所有子集组成的集合 \(\mathcal{P}(G)\) 上。如上所述，这定义了 \(G\) 在 \(\mathcal{P}(G)\) 上的一个群作用。注意若 \(S\) 是单元素集合 \(\{s\}\)，则 \(g \cdot S\) 是单元素集合 \(\{gsg^{-1}\}\)，所以 \(G\) 在自身所有子集上的这个作用可以看作 \(G\) 通过共轭作用在自身上的一种推广。

\begin{definition}
若存在 \(g \in G\) 使 \(T = gSg^{-1}\)（即当且仅当它们位于 \(G\) 通过共轭作用于其子集所得的同一轨道中），则称 \(G\) 的两个子集 \(S\) 和 \(T\) 在 \(G\) 中\textbf{共轭}。
\end{definition}

现在把命题 \ref{pro:4.2} 应用于 \(G\) 的共轭作用。命题 \ref{pro:4.2} 证明：若 \(S\) 是 \(G\) 的子集，则 \(S\) 的共轭的个数等于稳定子 \(G_S\) 的指数 \(|G : G_S|\). 对共轭作用，
\[
G_S = \{g \in G \mid gSg^{-1} = S\} = N_G(S)
\]
是 \(S\) 在 \(G\) 中的正规化子。我们把它总结为

\begin{proposition}\label{pro:4.6}
群 \(G\) 中子集 \(S\) 的共轭的个数是 \(S\) 的正规化子的指数 \(|G : N_G(S)|\). 特别地，\(G\) 中元素 \(s\) 的共轭的个数是 \(s\) 的中心化子的指数 \(|G : C_G(s)|\).
\end{proposition}

\begin{proof}
命题的第二个论断来自对 \(N_G(\{s\}) = C_G(s)\) 的观察。
\end{proof}

\(G\) 通过共轭作用在自身上把 \(G\) 划分成 \(G\) 的共轭类，其阶可以由命题 \ref{pro:4.6} 计算。由于这些共轭类的阶之和是 \(G\) 的阶，我们得到这些阶之间的如下重要关系。

\begin{theorem}[类方程]\label{thm:4.7}
设 \(G\) 是有限群，\(g_1, g_2, \ldots, g_r\) 是 \(G\) 的不含于中心 \(Z(G)\) 中的不同共轭类的代表元。则
\[
|G| = |Z(G)| + \sum_{i=1}^{r} |G : C_G(g_i)|.
\]
\end{theorem}

\begin{proof}
如上例 2 所述，单元素集合 \(\{x\}\) 是大小为 1 的共轭类当且仅当 \(x \in Z(G)\)，因为此时对所有 \(g \in G\) 有 \(gxg^{-1} = x\). 设 \(Z(G) = \{1, z_2, \ldots, z_m\}\)，设 \(K_1, K_2, \ldots, K_r\) 是 \(G\) 的不含于中心中的共轭类，并对每个 \(i\) 设 \(g_i\) 是 \(K_i\) 的代表元。则 \(G\) 的全部共轭类由
\[
\{1\}, \{z_2\}, \ldots, \{z_m\}, K_1, K_2, \ldots, K_r
\]
给出。由于它们划分 \(G\)，我们有
\[
|G| = \sum_{i=1}^{m} 1 + \sum_{i=1}^{r} |K_i| = |Z(G)| + \sum_{i=1}^{r} |G : C_G(g_i)|,
\]
其中 \(|K_i|\) 由命题 \ref{pro:4.6} 给出。这证明了类方程。
\end{proof}

特别要注意，类方程右端所有加项都是群阶的因子，因为它们是 \(G\) 的某个子群的指数。这限制了它们可能的取值（例如参见习题 6）。

\begin{example}
（1）类方程在阿贝尔群中不提供任何信息，因为共轭是平凡作用且所有共轭类的大小都是 1.

（2）在任意群 \(G\) 中都有 \(\langle g \rangle \leq C_G(g)\)；这个观察有助于尽量减少共轭类的计算。例如，在四元数群 \(Q_8\) 中我们看到 \(\langle i \rangle \leq C_{Q_8}(i) \leq Q_8\). 由于 \(i \notin Z(Q_8)\) 且 \(|Q_8 : \langle i \rangle| = 2\)，必有 \(C_{Q_8}(i) = \langle i \rangle\). 于是 \(i\) 在 \(Q_8\) 中恰有 2 个共轭，即 \(i\) 与 \(-i = kik^{-1}\). \(Q_8\) 中其他共轭类类似地确定，它们是
\[
\{1\},\quad \{-1\},\quad \{\pm i\},\quad \{\pm j\},\quad \{\pm k\}.
\]
前两个类构成 \(Z(Q_8)\)，这个群的类方程是
\[
|Q_8| = 2 + 2 + 2 + 2.
\]

（3）在 \(D_8\) 中我们也可以利用指数为 2 的三个子群都阿贝尔这一事实，迅速看出若 \(x \notin Z(D_8)\)，则 \(|C_{D_8}(x)| = 4\). \(D_8\) 的共轭类是
\[
\{1\},\quad \{r^2\},\quad \{r, r^3\},\quad \{s, sr^2\},\quad \{sr, sr^3\}.
\]
前两个类构成 \(Z(D_8)\)，这个群的类方程是
\[
|D_8| = 2 + 2 + 2 + 2.
\]
\end{example}

在讨论更多共轭的例子之前，我们给出类方程的两个重要推论。类方程的第一个应用是证明素数幂阶群有非平凡中心，这是研究素数幂阶群的起点（我们将在第 6 章回到这个话题）。

\begin{theorem}\label{thm:4.8}
若 \(p\) 是素数，\(P\) 是素数幂阶群 \(p^\alpha\)（\(\alpha \geq 1\)），则 \(P\) 有非平凡中心：\(Z(P) \neq 1\).
\end{theorem}

\begin{proof}
由类方程
\[
|P| = |Z(P)| + \sum_{i=1}^{r} |P : C_P(g_i)|,
\]
其中 \(g_1, \ldots, g_r\) 是不同非中心共轭类的代表元。按定义，对 \(i = 1, 2, \ldots, r\) 有 \(C_P(g_i) \neq P\)，故 \(p\) 整除 \(|P : C_P(g_i)|\). 由于 \(p\) 也整除 \(|P|\)，可推出 \(p\) 整除 \(|Z(P)|\)，因此中心非平凡。
\end{proof}

\begin{corollary}\label{cor:4.9}
若对某个素数 \(p\) 有 \(|P| = p^2\)，则 \(P\) 是阿贝尔群。更确切地说，\(P\) 同构于 \(Z_{p^2}\) 或 \(Z_p \times Z_p\).
\end{corollary}

\begin{proof}
由定理知 \(Z(P) \neq 1\)，故 \(P/Z(P)\) 是循环群。由第 3.1 节习题 36，\(P\) 是阿贝尔群。若 \(P\) 有 \(p^2\) 阶元素，则 \(P\) 是循环群。因此假设 \(P\) 的每个非单位元都是 \(p\) 阶的。设 \(x\) 是 \(P\) 的任意非单位元，\(y \in P - \langle x \rangle\). 由于 \(|\langle x, y \rangle| > |\langle x \rangle| = p\)，必有 \(P = \langle x, y \rangle\). \(x\) 和 \(y\) 的阶都是 \(p\)，故 \(\langle x \rangle \times \langle y \rangle = Z_p \times Z_p\). 由此直接得到映射 \((x^a, y^b) \mapsto x^a y^b\) 是从 \(\langle x \rangle \times \langle y \rangle\) 到 \(P\) 的同构。完成证明。
\end{proof}

\subsection*{\(S_n\) 中的共轭}

接下来考虑对称群中的共轭。熟悉线性代数的读者会认出，在矩阵群 \(GL_n(F)\) 中，共轭与“换基”是同一回事：\(A \mapsto PAP^{-1}\). \(S_n\) 中的情形是类似的：

\begin{proposition}\label{pro:4.10}
设 \(\sigma, \tau\) 是对称群 \(S_n\) 的元素，并假设 \(\sigma\) 有循环分解
\[
(a_1\, a_2\, \ldots\, a_{k_1})(b_1\, b_2\, \ldots\, b_{k_2}) \cdots .
\]
则 \(\tau\sigma\tau^{-1}\) 有循环分解
\[
(\tau(a_1)\, \tau(a_2)\, \ldots\, \tau(a_{k_1}))(\tau(b_1)\, \tau(b_2)\, \ldots\, \tau(b_{k_2})) \cdots,
\]
也就是说，\(\tau\sigma\tau^{-1}\) 由 \(\sigma\) 把循环分解中的每个条目 \(i\) 替换为条目 \(\tau(i)\) 得到。
\end{proposition}

\begin{proof}
注意若 \(\sigma(i) = j\)，则
\[
\tau\sigma\tau^{-1}(\tau(i)) = \tau(j).
\]
因此，若有序对 \(i, j\) 出现在 \(\sigma\) 的循环分解中，则有序对 \(\tau(i), \tau(j)\) 出现在 \(\tau\sigma\tau^{-1}\) 的循环分解中。完成证明。
\end{proof}

\begin{example}
设 \(\sigma = (1\,2)(3\,4\,5)(6\,7\,8\,9)\)，\(\tau = (1\,3\,5\,7)(2\,4\,6\,8)\). 则
\[
\tau\sigma\tau^{-1} = (3\,4)(5\,6\,7)(8\,1\,2\,9).
\]
\end{example}

\begin{definition}
（1）若 \(\sigma \in S_n\) 是长度分别为 \(n_1, n_2, \ldots, n_r\)（\(n_1 \geq n_2 \geq \cdots \geq n_r\)）的不交循环的乘积（包括其 1-循环），则称整数 \(n_1, n_2, \ldots, n_r\) 为 \(\sigma\) 的\textbf{循环型}。

（2）若 \(n \in \mathbb{Z}^+\)，则 \(n\) 的一个\textbf{分拆}是任意一个和为 \(n\) 的非递降正整数序列。
\end{definition}

注意由上一节的结果，置换的循环型是唯一的。例如，\(S_n\) 中一个 \(m\)-循环的循环型是 \(1, 1, \ldots, 1, m\)，其中 \(m\) 前面有 \(n-m\) 个 1.

\begin{proposition}\label{pro:4.11}
\(S_n\) 的两个元素在 \(S_n\) 中共轭，当且仅当它们有相同的循环型。\(S_n\) 的共轭类的个数等于 \(n\) 的分拆的个数。
\end{proposition}

\begin{proof}
由命题 \ref{pro:4.10}，共轭的置换有相同的循环型。反之，假设置换 \(\sigma_1\) 和 \(\sigma_2\) 有相同的循环型。按长度非递降地排列循环，包括 1-循环（若 \(\sigma_1\) 和 \(\sigma_2\) 有若干个相同长度的循环，则有若干种排法）。忽略括号，每个循环分解都是一个列表，其中 1 到 \(n\) 的所有整数恰好出现一次。定义 \(\tau\) 为把 \(\sigma_1\) 的列表中第 \(i\) 个整数映到 \(\sigma_2\) 的列表中第 \(i\) 个整数的函数。于是 \(\tau\) 是一个置换，且由于刻画循环分解的括号在每个列表中的位置相同，命题 \ref{pro:4.10} 保证 \(\tau\sigma_1\tau^{-1} = \sigma_2\)，因此 \(\sigma_1\) 和 \(\sigma_2\) 共轭。

由于 \(S_n\) 的共轭类与允许的循环型之间存在双射，且 \(S_n\) 中置换的每个循环型都是 \(n\) 的一个分拆，命题的第二个论断随之成立，完成证明。
\end{proof}

\begin{example}
（1）设 \(\sigma_1 = (1)(3\,5)(8\,9)(2\,4\,7\,6)\)，\(\sigma_2 = (3)(4\,7)(8\,1)(5\,2\,6\,9)\). 定义 \(\tau\) 为 \(\tau(1) = 3\)、\(\tau(3) = 4\)、\(\tau(5) = 7\)、\(\tau(8) = 8\) 等等。则
\[
\tau = (1\,3\,4\,2\,5\,7\,6\,9)(8)
\]
且 \(\tau\sigma_1\tau^{-1} = \sigma_2\).

（2）若在上一个例子中我们把 \(\sigma_2\) 重新排序为 \(\sigma_2 = (3)(8\,1)(4\,7)(5\,2\,6\,9)\)（交换两个长度为 2 的循环），则上面描述的相应 \(\tau\) 由 \(\tau(1) = 3\)、\(\tau(3) = 8\)、\(\tau(5) = 1\)、\(\tau(8) = 4\) 等等定义，得到置换
\[
\tau = (1\,3\,8\,4\,2\,5)(6\,9\,7),
\]
同样有 \(\tau\sigma_1\tau^{-1} = \sigma_2\)，这表明有许多元素把 \(\sigma_1\) 共轭为 \(\sigma_2\).

（3）若 \(n = 5\)，则 5 的分拆与相应的共轭类代表元（不写出 1-循环）如下表所示：

\begin{center}
\begin{tabular}{cc}
\hline
5 的分拆 & 共轭类的代表元 \\
\hline
\(1, 1, 1, 1, 1\) & \(1\) \\
\(1, 1, 1, 2\) & \((1\,2)\) \\
\(1, 1, 3\) & \((1\,2\,3)\) \\
\(1, 4\) & \((1\,2\,3\,4)\) \\
\(5\) & \((1\,2\,3\,4\,5)\) \\
\(1, 2, 2\) & \((1\,2)(3\,4)\) \\
\(2, 3\) & \((1\,2)(3\,4\,5)\) \\
\hline
\end{tabular}
\end{center}
\end{example}

命题 \ref{pro:4.11} 和命题 \ref{pro:4.6} 可以用来写出 \(S_n\) 中一些元素的中心化子。例如，若 \(\sigma\) 是 \(S_n\) 中的 \(m\)-循环，则 \(\sigma\) 的共轭的个数（即 \(m\)-循环的个数）是
\[
\frac{n \cdot (n-1) \cdots (n-m+1)}{m}.
\]
由命题 \ref{pro:4.6}，这是 \(\sigma\) 的中心化子的指数，故
\[
|C_{S_n}(\sigma)| = \frac{|S_n|}{n \cdot (n-1) \cdots (n-m+1)/m}.
\]
由于 \(|S_n| = n!\)，我们得到
\[
|C_{S_n}(\sigma)| = m \cdot (n-m)!.
\]
元素 \(\sigma\) 当然与 \(1, \sigma, \sigma^2, \ldots, \sigma^{m-1}\) 交换。它也与 \(S_n\) 中任何与 \(\sigma\) 不交的循环对应的置换交换，这样的置换有 \((n-m)!\) 个（在不出现于 \(\sigma\) 中的数上的完整对称群）。这两类元素的乘积已经给出了 \(m \cdot (n-m)!\) 个与 \(\sigma\) 交换的元素。由上面的阶的计算，这就是 \(\sigma\) 在 \(S_n\) 中的全部中心化子。具体地说，若 \(\sigma\) 是 \(S_n\) 中的 \(m\)-循环，则
\[
C_{S_n}(\sigma) = \{\sigma^i \tau \mid 0 \leq i \leq m-1,\ \tau \in S_{n-m}\},
\]
其中 \(S_{n-m}\) 表示 \(S_n\) 中固定 \(m\)-循环 \(\sigma\) 中出现的所有整数的子群（若 \(m = n\) 或 \(m = n-1\)，则为单位子群）。

例如，\(\sigma = (1\,3\,5)\) 在 \(S_7\) 中的中心化子是由
\[
\{(1\,3\,5)^i \tau \mid i = 0, 1 \text{ 或 } 2,\ \text{且 } \tau \text{ 固定 } 1, 3, 5\}
\]
给出的子群。注意 \(\tau \in S_A\)，其中 \(A = \{2, 4, 6, 7\}\)，所以 \(\tau\) 有 \(4!\) 种选择，中心化子的阶为 \(3 \cdot 4! = 72\).

我们将在后面的习题和第 5 章讨论 \(S_n\) 中其他元素的中心化子。

我们可以用关于 \(S_n\) 共轭类的这一讨论来给出 \(A_5\) 单性的一个组合证明。我们首先注意到，群 \(G\) 的正规子群是 \(G\) 的共轭类的并，即
\begin{center}
若 \(H \trianglelefteq G\)，则对 \(G\) 的每个共轭类 \(K\)，或者 \(K \subseteq H\)，或者 \(K \cap H = \varnothing\).
\end{center}
这是因为若 \(x \in K \cap H\)，则对所有 \(g \in G\) 有 \(gxg^{-1} \in gHg^{-1}\). 由于 \(H\) 正规，\(gHg^{-1} = H\)，故 \(H\) 包含 \(x\) 的所有共轭，即 \(K \subseteq H\).

\begin{theorem}\label{thm:4.12}
\(A_5\) 是单群。
\end{theorem}

\begin{proof}
我们首先求出 \(A_5\) 的共轭类及其大小。命题 \ref{pro:4.11} 不能直接应用，因为相同循环型的两个元素（它们在 \(S_5\) 中共轭）未必在 \(A_5\) 中共轭。习题 19 至 22 详细分析 \(S_n\) 中的类与 \(A_n\) 中的类的关系。

我们已经看到，偶置换的循环型的代表元可以取为
\[
1,\quad (1\,2\,3),\quad (1\,2\,3\,4\,5),\quad (1\,2)(3\,4).
\]
3-循环和 5-循环在 \(S_5\) 中的中心化子已在上面确定，检查其中哪些元素包含在 \(A_5\) 中，我们看到
\[
C_{A_5}((1\,2\,3)) = \langle (1\,2\,3) \rangle \quad \text{和} \quad C_{A_5}((1\,2\,3\,4\,5)) = \langle (1\,2\,3\,4\,5) \rangle.
\]
这些群的阶分别为 3 和 5（指数分别为 20 和 12），所以在 \(A_5\) 中 \((1\,2\,3)\) 有 20 个不同的共轭，\((1\,2\,3\,4\,5)\) 有 12 个不同的共轭。由于 \(S_5\) 中总共有二十个 3-循环（第 1.3 节习题 16），且它们都落在 \(A_5\) 中，我们看到
\begin{center}
所有二十个 3-循环在 \(A_5\) 中共轭。
\end{center}
\(A_5\) 中总共有二十四个 5-循环，但 5-循环 \((1\,2\,3\,4\,5)\) 只有 12 个不同的共轭。因此某个 5-循环 \(\sigma\) 在 \(A_5\) 中不与 \((1\,2\,3\,4\,5)\) 共轭（事实上，\((1\,3\,5\,2\,4)\) 在 \(A_5\) 中不与 \((1\,2\,3\,4\,5)\) 共轭，因为命题 \ref{pro:4.11} 证明方法表明把 \((1\,2\,3\,4\,5)\) 共轭为 \((1\,3\,5\,2\,4)\) 的任何 \(S_5\) 中元素必是奇置换）。如上所述，我们看到 \(\sigma\) 在 \(A_5\) 中也有 12 个不同的共轭，因此 5-循环在 \(A_5\) 中落在两个共轭类中，每个有 12 个元素。

由于 3-循环和 5-循环已经包含了所有奇阶非单位元，\(A_5\) 余下的 15 个非单位元必是 2 阶的，因而循环型为 \((2,2)\). 容易看出 \((1\,2)(3\,4)\) 与 \((1\,3)(2\,4)\) 交换，但不与 \(A_5\) 中任何奇阶元素交换。由此推出 \(|C_{A_5}((1\,2)(3\,4))| = 4\). 于是 \((1\,2)(3\,4)\) 在 \(A_5\) 中有 15 个不同的共轭，因此
\begin{center}
\(A_5\) 中所有 15 个 2 阶元素都与 \((1\,2)(3\,4)\) 共轭。
\end{center}
总之，\(A_5\) 的共轭类的阶为 1, 15, 20, 12 和 12.

现在假设 \(H\) 是 \(A_5\) 的正规子群。那么如我们上面注意到的，\(H\) 会是 \(A_5\) 的若干个共轭类的并。于是 \(H\) 的阶既是 60（\(A_5\) 的阶）的因子，又是整数集合 \(\{1, 12, 12, 15, 20\}\)（\(A_5\) 中各共轭类的大小）中若干个的和。稍加检验可知唯一的可能是 \(|H| = 1\) 或 \(|H| = 60\)，故 \(A_5\) 没有非平凡正规子群。
\end{proof}

\subsection*{右群作用}

如第 1.7 节所述，在作用的定义中群元素出现在集合元素的左边，所以我们的作用概念更确切地说可以称为\textbf{左群作用}。人们可以类似地定义群 \(G\) 在非空集合 \(A\) 上的\textbf{右群作用}为从 \(A \times G\) 到 \(A\) 的映射，记作 \(a \cdot g\)（\(a \in A\)，\(g \in G\)），满足公理：
\begin{enumerate}[label=(\arabic*)]
\item \((a \cdot g_1) \cdot g_2 = a \cdot (g_1 g_2)\) 对所有 \(a \in A\) 和 \(g_1, g_2 \in G\)；
\item \(a \cdot 1 = a\) 对所有 \(a \in A\).
\end{enumerate}

在群论的许多文献中，共轭被写成右群作用，使用如下记号：
\[
a^g = g^{-1}ag \quad \text{对所有 } g, a \in G.
\]
类似地，对 \(G\) 的子集 \(S\) 定义 \(S^g = g^{-1}Sg\). 用这个记号，右作用的两条公理验证如下：
\[
(a^{g_1})^{g_2} = g_2^{-1}(g_1^{-1}ag_1)g_2 = (g_1g_2)^{-1}a(g_1g_2) = a^{g_1g_2}
\]
和
\[
a^1 = 1^{-1}a1 = a
\]
对所有 \(g_1, g_2, a \in G\). 于是这个群在自身上右作用的两条公理呈现为熟悉的“指数律”形式。（注意群元素 \(a\) 的整数幂 \(a^n\) 与 \(a\) 的共轭 \(a^g\) 容易通过指数的性质区分：\(n \in \mathbb{Z}\) 但 \(g \in G\)。）由于共轭在群论中无处不在，这个记号是一种有用而高效的简写（不必总是写 \(gag^{-1}\) 或 \(g \cdot a\) 来表示左共轭作用）。

对任意的群作用，容易验证：若给定 \(G\) 在 \(A\) 上的左群作用，则由 \(a \cdot g = g^{-1} \cdot a\) 定义的映射 \(A \times G \to A\) 是右群作用。反之，给定 \(G\) 在 \(A\) 上的右群作用，我们可以通过 \(g \cdot a = a \cdot g^{-1}\) 得到一个左群作用。称这些为\textbf{相应的群作用}。换句话说，对相应的群作用，\(g\) 从左边作用与 \(g^{-1}\) 从右边作用的方式相同。这对共轭作用尤其明显，因为“\(a\) 被 \(g\) 左共轭”即 \(gag^{-1}\)，与“\(a\) 被 \(g^{-1}\) 右共轭”即 \(a^{g^{-1}}\) 是同一个群元素。因此群的两个元素或子集“左共轭”当且仅当它们“右共轭”，所以“共轭”这一关系对左右相应的作用是相同的。更一般地，容易验证（习题 1）：对任何相应的左右作用，轨道是相同的。

我们一贯使用左作用，因为它们与在左边施用函数的记号（即记号 \(\varphi(g)\)）相容；这样，在子群的左陪集上的左乘就是一个左作用。类似地，在子群的右陪集上的右乘是一个右作用，只要函数 \(\varphi : G \to S_A\) 写在右边（即 \((g_1g_2)\varphi = (g_1\varphi)(g_2\varphi)\)，并且 \(S_A\) 中的置换作为从 \(A\) 到自身的函数写在右边），相应的置换表示就是同态。有些情形下一个集合容许群 \(G\) 的两个作用：一个自然地在左边，另一个在右边，因此熟悉这两类作用是有用的。

\subsection*{习题}
设 \(G\) 是群。

\begin{enumerate}
\item 假设 \(G\) 左作用在集合 \(A\) 上，记作 \(g \cdot a\)（对所有 \(g \in G\)、\(a \in A\)）。用 \(a \cdot g\) 表示相应的右作用。证明由
\[
a \sim b \quad \text{当且仅当} \quad \text{对某个 } g \in G \text{ 有 } a = g \cdot b
\]
和
\[
a \sim' b \quad \text{当且仅当} \quad \text{对某个 } g \in G \text{ 有 } a = b \cdot g
\]
定义的（等价）关系 \(\sim\) 和 \(\sim'\) 是同一个关系（即 \(a \sim b\) 当且仅当 \(a \sim' b\)）。

\item 求下列群的所有共轭类及其大小：

（a）\(D_8\) \qquad（b）\(Q_8\) \qquad（c）\(A_4\).

\item 求下列群的所有共轭类及其大小：

（a）\(Z_2 \times S_3\) \qquad（b）\(S_3 \times S_3\) \qquad（c）\(Z_3 \times A_4\).

\item 证明：若 \(S \leq G\) 且 \(g \in G\)，则 \(gN_G(S)g^{-1} = N_G(gSg^{-1})\) 且 \(gC_G(S)g^{-1} = C_G(gSg^{-1})\).

\item 若 \(G\) 的中心指数为 \(n\)，证明每个共轭类最多有 \(n\) 个元素。

\item 假设 \(G\) 是 15 阶非阿贝尔群。证明 \(Z(G) = 1\). 利用对所有 \(g \in G\) 有 \(\langle g \rangle \leq C_G(g)\) 这一事实证明 \(G\) 至多有一种可能的类方程。〔使用第 3.1 节习题 36。〕

\item 对 \(n = 3, 4, 6, 7\) 列出 \(n\) 的分拆，并给出 \(S_n\) 相应共轭类的代表元。

\item 证明对所有 \(n \geq 3\) 有 \(Z(S_n) = 1\).

\item 证明对所有 \(n \geq 4\) 有 \(|C_{S_n}((1\,2)(3\,4))| = 8 \cdot (n-4)!\). 具体确定这个中心化子中的元素。

\item 设 \(\sigma\) 是 \(S_5\) 中的 5-循环 \((1\,2\,3\,4\,5)\). 在 (a) 到 (c) 中各求一个具体的 \(\tau \in S_5\) 实现指定的共轭：

（a）\(\tau\sigma\tau^{-1} = \sigma^2\)；

（b）\(\tau\sigma\tau^{-1} = \sigma^{-1}\)；

（c）\(\tau\sigma\tau^{-1} = \sigma^{-2}\).

\item 在 (a)–(d) 中各判断 \(\sigma_1\) 与 \(\sigma_2\) 是否共轭。若共轭，给出一个具体的置换 \(\tau\) 使 \(\tau\sigma_1\tau^{-1} = \sigma_2\).

（a）\(\sigma_1 = (1\,2)(3\,4\,5)\) 和 \(\sigma_2 = (1\,2\,3)(4\,5)\)；

（b）\(\sigma_1 = (1\,5)(3\,7\,2)(10\,6\,8\,11)\) 和 \(\sigma_2 = (3\,7\,5\,10)(4\,9)(13\,11\,2)\)；

（c）\(\sigma_1 = (1\,5)(3\,7\,2)(10\,6\,8\,11)\) 和 \(\sigma_2 = \sigma_1^2\)；

（d）\(\sigma_1 = (1\,3)(2\,4\,6)\) 和 \(\sigma_2 = (3\,5)(2\,4)(5\,6)\).

\item 求 \(S_5\) 和 \(S_{12}\) 中每个 4 阶元素共轭类的一个代表元。

\item 求所有恰有两个共轭类的有限群。

\item 在第 2 节习题 1 中，对克莱因四元群 \(V\) 的元素 \(\{1, a, b, c\}\) 选了两种标号，给出 \(V\) 到 \(S_4\) 的左正则表示的两种版本。设 \(\pi_1\) 是该习题 (a) 中得到的正则表示版本，\(\pi_2\) 是 (b) 中标号得到的版本。设 \(\tau = (2\,4)\). 证明对每个 \(g \in V\) 有 \(\tau \circ \pi_1(g) \circ \tau^{-1} = \pi_2(g)\)（即用 \(\tau\) 共轭把 \(\pi_1\) 的像逐元素地送到 \(\pi_2\) 的像）。

\item 求 \(S_8\) 的一个元素，它把该节习题 3 的 (a) 中得到的 \(S_8\) 的子群共轭为同一习题 (b) 中得到的 \(S_8\) 的子群（这两个子群都同构于 \(D_8\)）。

\item 求 \(S_4\) 的一个元素，它把该节习题 5 的 (a) 中得到的 \(S_4\) 的子群共轭为同一习题 (b) 中得到的 \(S_4\) 的子群（这两个子群都同构于 \(D_8\)）。

\item 设 \(A\) 是非空集合，\(X\) 是 \(S_A\) 的任意子集。令
\[
F(X) = \{a \in A \mid \sigma(a) = a \text{ 对所有 } \sigma \in X\}
\]
为 \(X\) 的\textbf{不动点集}。令 \(M(X) = A - F(X)\) 为被 \(X\) 中某个元素移动的元素。令 \(D = \{\sigma \in S_A \mid |M(\sigma)| < \infty\}\). 证明 \(D\) 是 \(S_A\) 的正规子群。

\item 设 \(A\) 是集合，\(H\) 是 \(S_A\) 的子群，\(F(H)\) 是上一习题定义的 \(H\) 在 \(A\) 上的不动点。证明：若 \(\tau \in N_{S_A}(H)\)，则 \(\tau\) 稳定集合 \(F(H)\) 及其补集 \(A - F(H)\).

\item 假设 \(H\) 是 \(G\) 的正规子群，\(K\) 是 \(G\) 的含于 \(H\) 中的共轭类，\(x \in K\). 证明 \(K\) 是 \(H\) 中 \(k\) 个大小相等的共轭类的并，其中 \(k = |G : HC_G(x)|\). 由此推出：\(S_n\) 中由偶置换组成的共轭类在 \(A_n\) 的作用下或者是单个共轭类，或者是两个大小相同的类之并。〔令 \(A = C_G(x)\)、\(B = H\)，故 \(A \cap B = C_H(x)\). 画出与第二同构定理相关的格图，并解释相应的指数。另见第 1 节习题 9。〕

\item 设 \(\sigma \in A_n\). 证明 \(\sigma\) 在 \(S_n\) 中的共轭类（即所有与 \(\sigma\) 有相同循环型的元素）中的所有元素在 \(A_n\) 中共轭，当且仅当 \(\sigma\) 与一个奇置换交换。〔使用前一习题。〕

\item 设 \(K\) 是 \(S_n\) 中的共轭类，并假设 \(K \subseteq A_n\). 证明：\(\sigma \in S_n\) 不与任何奇置换交换，当且仅当 \(\sigma\) 的循环型由互不相同的奇数组成。由此推出：\(K\) 在 \(A_n\) 中由两个共轭类组成，当且仅当 \(K\) 中元素的循环型由互不相同的奇数组成。〔首先假设 \(\sigma \in K\) 不与任何奇置换交换。注意 \(\sigma\) 与其循环分解中各个单独的循环交换——用它证明其所有循环必为奇数长度。若有两个循环长度相同（设为 \(k\)），找出交换它们且与 \(\sigma\) 交换的 \(k\) 个对换之积。反之，若 \(\sigma\) 的循环型由互不相同的整数组成，证明 \(\sigma\) 只与其循环分解中循环生成的群交换。〕

\item 证明：若 \(n\) 是奇数，则所有 \(n\)-循环的集合在 \(A_n\) 中由两个大小相等的共轭类组成。

\item 回忆（参见第 2.4 节习题 16），\(G\) 的真子群 \(M\) 称为\textbf{极大子群}，若只要 \(M \leq H \leq G\)，就有 \(H = M\) 或 \(H = G\). 证明：若 \(M\) 是 \(G\) 的极大子群，则 \(N_G(M) = M\) 或 \(N_G(M) = G\). 由此推出：若 \(M\) 是 \(G\) 的不正规极大子群，则包含在 \(M\) 的共轭中的 \(G\) 的非单位元最多有 \((|M| - 1)|G : M|\) 个。

\item 假设 \(H\) 是有限群 \(G\) 的真子群。证明 \(G \neq \bigcup_{g \in G} gHg^{-1}\)，即 \(G\) 不是任何真子群的共轭之并。〔把 \(H\) 放入某个极大子群并使用前一习题。〕

\item 设 \(G = GL_2(\mathbb{C})\)，\(H = \left\{ \begin{pmatrix} a & b \\ 0 & c \end{pmatrix} \ \middle|\ a, b, c \in \mathbb{C},\ ac \neq 0 \right\}\). 证明 \(G\) 的每个元素都共轭于子群 \(H\) 的某个元素，并由此推出 \(G\) 是 \(H\) 的共轭之并。〔证明 \(GL_2(\mathbb{C})\) 的每个元素都有特征向量。〕

\item 设 \(G\) 是有限集合 \(A\) 上的传递置换群，\(|A| > 1\). 证明存在 \(\sigma \in G\) 使得对所有 \(a \in A\) 有 \(\sigma(a) \neq a\)（这样的元素 \(\sigma\) 称为\textbf{无不动点}的）。

\item 设 \(g_1, g_2, \ldots, g_r\) 是有限群 \(G\) 的共轭类的代表元，并假设这些元素两两交换。证明 \(G\) 是阿贝尔群。

\item 设 \(p\) 和 \(q\) 是素数且 \(p < q\). 证明 \(pq\) 阶非阿贝尔群 \(G\) 有一个指数为 \(q\) 的非正规子群，从而存在到 \(S_q\) 的单同态。由此推出 \(G\) 同构于 \(S_q\) 中由 \(q\)-循环 \((1\,2\,\ldots\,q)\) 生成的循环群的正规化子的某个子群。

\item 设 \(p\) 是素数，\(G\) 是 \(p^\alpha\) 阶群。证明对每个 \(0 \leq \beta \leq \alpha\)，\(G\) 都有 \(p^\beta\) 阶子群。〔使用定理 \ref{thm:4.8} 并对 \(\alpha\) 作归纳。〕

\item 若 \(G\) 是奇数阶群，证明对任意非单位元 \(x \in G\)，\(x\) 与 \(x^{-1}\) 在 \(G\) 中不共轭。

\item 用二面体群 \(D_{2n}\) 的通常生成元与关系（参见第 1.2 节）证明：当 \(n = 2k\) 为偶数时，\(D_{2n}\) 的共轭类如下：\(\{1\}\)、\(\{r^k\}\)、\(\{r^{\pm 1}\}\)、\(\{r^{\pm 2}\}\)、\(\ldots\)、\(\{r^{\pm(k-1)}\}\)、\(\{sr^{2b} \mid b = 1, \ldots, k\}\) 和 \(\{sr^{2b-1} \mid b = 1, \ldots, k\}\). 写出 \(D_{2n}\) 的类方程。

\item 对 \(n = 2k + 1\) 为奇数，证明 \(D_{2n}\) 的共轭类是 \(\{1\}\)、\(\{r^{\pm 1}\}\)、\(\{r^{\pm 2}\}\)、\(\ldots\)、\(\{r^{\pm k}\}\)、\(\{sr^b \mid b = 1, \ldots, n\}\). 写出 \(D_{2n}\) 的类方程。

\item 这个习题给出 \(S_n\) 中每个共轭类大小的公式。设 \(\sigma\) 是 \(S_n\) 中的置换，\(m_1, m_2, \ldots, m_s\) 是 \(\sigma\) 的循环型中出现的不同整数（包括 1-循环）。对每个 \(i \in \{1, 2, \ldots, s\}\) 假设 \(\sigma\) 有 \(k_i\) 个长度为 \(m_i\) 的循环（故 \(\sum_{i=1}^{s} k_i m_i = n\)）。证明 \(\sigma\) 的共轭个数为
\[
\frac{n!}{(k_1!\, m_1^{k_1})(k_2!\, m_2^{k_2}) \cdots (k_s!\, m_s^{k_s})}.
\]
〔参见第 1.3 节习题 6 和 7，那里在一些特殊情形给出了这个公式。〕

\item 证明：若 \(p\) 是素数，\(P\) 是 \(S_p\) 的 \(p\) 阶子群，则 \(|N_{S_p}(P)| = p(p-1)\).〔论证 \(P\) 的每个共轭恰含 \(p-1\) 个 \(p\)-循环，并用 \(p\)-循环个数的公式计算 \(N_{S_p}(P)\) 在 \(S_p\) 中的指数。〕

\item 设 \(p\) 是素数。用取整函数给出 \(S_n\) 中 \(p\) 阶元素的共轭类个数的公式。

\item 设 \(\pi : G \to S_G\) 是 \(G\) 通过左乘作用在自身上所提供的左正则表示。对每个 \(g \in G\) 把置换 \(\pi(g)\) 记作 \(\sigma_g\)，故对所有 \(x \in G\) 有 \(\sigma_g(x) = gx\). 设 \(\lambda : G \to S_G\) 是相应的 \(G\) 在自身上右作用所提供的置换表示，对每个 \(h \in G\) 把置换 \(\lambda(h)\) 记作 \(\tau_h\). 于是对所有 \(x \in G\) 有 \(\tau_h(x) = xh^{-1}\)（\(\lambda\) 称为 \(G\) 的\textbf{右正则表示}）。

（a）证明对所有 \(g, h \in G\)，\(\sigma_g\) 与 \(\tau_h\) 交换。（于是 \(\pi(G)\) 在 \(S_G\) 中的中心化子包含同构于 \(G\) 的子群 \(\lambda(G)\)。）

（b）证明 \(\sigma_g = \tau_g\) 当且仅当 \(g\) 是 \(G\) 的中心中阶为 1 或 2 的元素。

（c）证明 \(\sigma_g = \tau_h\) 当且仅当 \(g\) 和 \(h\) 都在 \(G\) 的中心中。由此推出 \(\pi(G) \cap \lambda(G) = \pi(Z(G)) = \lambda(Z(G))\).
\end{enumerate}

\section{自同构}

\begin{definition}
设 \(G\) 是群。从 \(G\) 到自身的同构称为 \(G\) 的\textbf{自同构}。\(G\) 的所有自同构的集合记作 \(\operatorname{Aut}(G)\).
\end{definition}

我们把 \(\operatorname{Aut}(G)\) 在自同构的复合下构成群（\(G\) 的\textbf{自同构群}）这一简单验证留作练习（自同构的复合有定义，因为每个自同构的定义域与值域相同）。注意群 \(G\) 的自同构特别是集合 \(G\) 的置换，所以 \(\operatorname{Aut}(G)\) 是 \(S_G\) 的子群。

群 \(G\) 的自同构最重要的例子之一由被 \(G\) 中一个固定元素共轭给出。下一个结果在稍微更一般的框架中讨论它。

\begin{proposition}\label{pro:4.13}
设 \(H\) 是群 \(G\) 的正规子群。则 \(G\) 通过共轭作为 \(H\) 的自同构作用在 \(H\) 上。更具体地说，\(G\) 通过共轭在 \(H\) 上的作用对每个 \(g \in G\) 定义为
\[
h \mapsto ghg^{-1} \quad \text{对每个 } h \in H.
\]
对每个 \(g \in G\)，被 \(g\) 共轭是 \(H\) 的一个自同构。这个作用所提供的置换表示是 \(G\) 到 \(\operatorname{Aut}(H)\) 的同态，其核为 \(C_G(H)\). 特别地，\(G/C_G(H)\) 同构于 \(\operatorname{Aut}(H)\) 的一个子群。
\end{proposition}

\begin{proof}
（参见第 1.7 节习题 17。）设 \(\varphi_g\) 表示被 \(g\) 共轭。注意由于 \(g\) 正规化 \(H\)，\(\varphi_g\) 把 \(H\) 映到自身。由于我们已经看到共轭定义一个作用，故 \(\varphi_1 = \mathrm{id}\)（\(H\) 上的恒等映射）且对所有 \(a, b \in G\) 有 \(\varphi_a \circ \varphi_b = \varphi_{ab}\). 因此每个 \(\varphi_g\) 给出 \(H\) 到自身的一个双射，因为它有双边逆 \(\varphi_{g^{-1}}\). 每个 \(\varphi_g\) 是从 \(H\) 到 \(H\) 的同态，因为对所有 \(h, k \in H\)，
\[
\varphi_g(hk) = g(hk)g^{-1} = gh(gg^{-1})kg^{-1} = (ghg^{-1})(gkg^{-1}) = \varphi_g(h)\varphi_g(k).
\]
这证明了被 \(G\) 中任意固定元素共轭定义 \(H\) 的一个自同构。

由上述评注，由 \(\psi(g) = \varphi_g\) 定义的置换表示 \(\psi : G \to S_H\)（我们已经证明它是同态）的像包含在子群 \(\operatorname{Aut}(H)\) 中。最后，
\[
\ker \psi = \{g \in G \mid \varphi_g = \mathrm{id}\} = \{g \in G \mid ghg^{-1} = h \text{ 对所有 } h \in H\} = C_G(H).
\]
第一同构定理蕴含命题的最后一个论断。
\end{proof}

命题 \ref{pro:4.13} 表明群通过共轭作为保结构的置换、即自同构作用在正规子群上。特别地，这个作用必把子群送到子群，把 \(n\) 阶元素送到 \(n\) 阶元素，等等。这个命题的两个具体应用在下面两个推论中描述。

\begin{corollary}\label{cor:4.14}
若 \(K\) 是群 \(G\) 的任意子群且 \(g \in G\)，则 \(K \cong gKg^{-1}\). 共轭的元素和共轭的子群有相同的阶。
\end{corollary}

\begin{proof}
在命题中取 \(G = H\) 表明被 \(g \in G\) 共轭是 \(G\) 的自同构，推论随之成立。
\end{proof}

\begin{corollary}\label{cor:4.15}
对群 \(G\) 的任意子群 \(H\)，商群 \(N_G(H)/C_G(H)\) 同构于 \(\operatorname{Aut}(H)\) 的一个子群。特别地，\(G/Z(G)\) 同构于 \(\operatorname{Aut}(G)\) 的一个子群。
\end{corollary}

\begin{proof}
由于 \(H\) 是群 \(N_G(H)\) 的正规子群，把命题 \ref{pro:4.13} 应用于 \(N_G(H)\)（让它扮演 \(G\) 的角色）即得第一个论断。第二个论断是 \(H = G\) 的特殊情形，此时 \(N_G(G) = G\) 且 \(C_G(G) = Z(G)\).
\end{proof}

\begin{definition}
设 \(G\) 是群，\(g \in G\). 被 \(g\) 共轭称为 \(G\) 的\textbf{内自同构}，\(\operatorname{Aut}(G)\) 中由所有内自同构组成的子群记作 \(\operatorname{Inn}(G)\).
\end{definition}

注意 \(G\) 的内自同构的集合实际上是 \(\operatorname{Aut}(G)\) 的子群，且由推论 \ref{cor:4.15} 有 \(\operatorname{Inn}(G) \cong G/Z(G)\). 还要注意，若 \(H\) 是 \(G\) 的正规子群，则 \(G\) 的元素在限制到 \(H\) 时的共轭是 \(H\) 的自同构，但不必是 \(H\) 的内自同构（我们将看到这一点）。

\begin{example}
（1）群 \(G\) 是阿贝尔群当且仅当每个内自同构都是平凡的。若 \(H\) 是 \(G\) 的阿贝尔正规子群且 \(H\) 不含于 \(Z(G)\) 中，则存在某个 \(g \in G\)，使被 \(g\) 共轭限制到 \(H\) 上不是 \(H\) 的内自同构。一个具体例子是 \(G = A_4\)，\(H\) 是 \(G\) 中的克莱因四元群，\(g\) 是任意 3-循环。

（2）由于 \(Z(Q_8) = \langle -1 \rangle\)，我们有 \(\operatorname{Inn}(Q_8) \cong V_4\).

（3）由于 \(Z(D_8) = \langle r^2 \rangle\)，我们有 \(\operatorname{Inn}(D_8) \cong V_4\).

（4）由于对所有 \(n \geq 3\) 有 \(Z(S_n) = 1\)，我们有 \(\operatorname{Inn}(S_n) \cong S_n\).
\end{example}

推论 \ref{cor:4.15} 表明，我们拥有的关于群 \(G\) 的子群 \(H\) 的自同构群的任何信息都可以转化为关于 \(N_G(H)/C_G(H)\) 的信息。例如，若 \(H \cong Z_2\)，则由于 \(H\) 有唯一的 1 阶和 2 阶元素，推论 \ref{cor:4.14} 迫使 \(\operatorname{Aut}(H) = 1\). 于是若 \(H \cong Z_2\)，则 \(N_G(H) = C_G(H)\)；若此外 \(H\) 是 \(G\) 的正规子群，则 \(H \leq Z(G)\)（参见第 2.2 节习题 10）。

尽管上面的例子相当平凡，它表明 \(G\) 通过共轭作用在正规子群 \(H\) 上可以受关于 \(H\) 的自同构群的知识所限制。这反过来可以用来研究 \(G\) 的结构，并将在我们在 5.5 节考虑半直积时导向一些分类定理。

一个将在后面各节使用的概念最自然地在此时引入：

\begin{definition}
群 \(G\) 的子群 \(H\) 称为在 \(G\) 中\textbf{特征的}，记作 \(H\ \mathrm{char}\ G\)，若 \(G\) 的每个自同构都把 \(H\) 映到自身，即对所有 \(\sigma \in \operatorname{Aut}(G)\) 有 \(\sigma(H) = H\).
\end{definition}

我们将要在后面使用的关于特征子群的结果（其证明留作习题）是：

（1）特征子群是正规的；

（2）若 \(H\) 是 \(G\) 中给定阶的唯一子群，则 \(H\) 在 \(G\) 中特征；

（3）若 \(K\ \mathrm{char}\ H\) 且 \(H \trianglelefteq G\)，则 \(K \trianglelefteq G\)（所以，尽管“正规性”不是传递性质（即正规子群的正规子群不必正规），正规子群的特征子群是正规的）。

因此我们可以把特征子群看作“强正规”子群。例如，性质 (2) 和定理 2.7 蕴含循环群的每个子群都是特征的。

本节最后给出关于具体群的自同构群的一些结果。

\begin{proposition}\label{pro:4.16}
\(n\) 阶循环群的自同构群同构于 \((\mathbb{Z}/n\mathbb{Z})^\times\)，一个 \(\varphi(n)\) 阶阿贝尔群（其中 \(\varphi\) 是欧拉函数）。
\end{proposition}

\begin{proof}
设 \(x\) 是循环群 \(Z_n\) 的生成元。若 \(\psi \in \operatorname{Aut}(Z_n)\)，则对某个 \(a \in \mathbb{Z}\) 有 \(\psi(x) = x^a\)，且整数 \(a\) 唯一确定 \(\psi\). 把这个自同构记作 \(\psi_a\). 通常地，由于 \(|x| = n\)，整数 \(a\) 只在模 \(n\) 的意义下确定。由于 \(\psi_a\) 是自同构，\(x\) 与 \(x^a\) 必有相同的阶，故 \((a, n) = 1\). 此外，对每个与 \(n\) 互素的 \(a\)，映射 \(x \mapsto x^a\) 是 \(Z_n\) 的自同构。因此我们有一个满射
\[
\Psi : \operatorname{Aut}(Z_n) \to (\mathbb{Z}/n\mathbb{Z})^\times, \quad \psi_a \mapsto a \pmod{n}.
\]
映射 \(\Psi\) 是同态，因为对所有 \(\psi_a, \psi_b \in \operatorname{Aut}(Z_n)\)，
\[
\psi_a \circ \psi_b(x) = \psi_a(x^b) = (x^b)^a = x^{ab} = \psi_{ab}(x),
\]
故
\[
\Psi(\psi_a \circ \psi_b) = \Psi(\psi_{ab}) = ab \pmod{n} = \Psi(\psi_a)\Psi(\psi_b).
\]
最后，\(\Psi\) 显然是单射，因此是同构。
\end{proof}

\(\operatorname{Aut}(Z_n)\) 的同构类型的完整描述在 9.5 节末给出。

\begin{example}
假设 \(G\) 是 \(pq\) 阶群，其中 \(p\) 和 \(q\) 是素数（不必不同）且 \(p \leq q\). 若 \(p \nmid q - 1\)，我们证明 \(G\) 是阿贝尔群。

若 \(Z(G) \neq 1\)，拉格朗日定理迫使 \(G/Z(G)\) 是循环群，故由第 3.1 节习题 36，\(G\) 是阿贝尔群。因此可设 \(Z(G) = 1\).

若 \(G\) 的每个非单位元的阶都是 \(p\)，则每个非单位元的中心化子指数为 \(q\)，故 \(G\) 的类方程为
\[
pq = 1 + kq.
\]
这不可能，因为 \(q\) 整除 \(pq\) 和 \(kq\)，但不整除 1. 因此 \(G\) 含有 \(q\) 阶元素 \(x\).

设 \(H = \langle x \rangle\). 由于 \(H\) 的指数为 \(p\)，而 \(p\) 是整除 \(|G|\) 的最小素数，由推论 \ref{cor:4.5} 子群 \(H\) 在 \(G\) 中正规。由于 \(Z(G) = 1\)，必有 \(C_G(H) = H\). 于是由推论 \ref{cor:4.15}，\(G/H = N_G(H)/C_G(H)\) 是 \(p\) 阶群，同构于 \(\operatorname{Aut}(H)\) 的一个子群。但由命题 \ref{pro:4.16}，\(\operatorname{Aut}(H)\) 的阶为 \(\varphi(q) = q - 1\)，由拉格朗日定理这将蕴含 \(p \mid q - 1\)，与假设矛盾。这表明 \(G\) 必是阿贝尔群。
\end{example}

可以验证，\(pq\) 阶群（其中 \(p, q\) 是不同素数，\(p < q\) 且 \(p \nmid q - 1\)）都是循环群（见习题）。这是第一个存在唯一同构类型的合数阶群的情形。例如，每个 15 阶群都是循环群。

下一个命题总结了关于已知群的自同构群的一些结果，将在后面证明。这个命题的第 3 部分说明了向量空间理论如何在群论中发挥作用。

\begin{proposition}\label{pro:4.17}
\begin{enumerate}[label=(\arabic*)]
\item 若 \(p\) 是奇素数且 \(n \in \mathbb{Z}^+\)，则 \(p\) 阶循环群的自同构群是 \(p-1\) 阶循环群。更一般地，\(p^n\) 阶循环群的自同构群是 \(p^{n-1}(p-1)\) 阶循环群（参见 9.5 节推论 20）。

\item 对所有 \(n \geq 3\)，\(2^n\) 阶循环群的自同构群同构于 \(Z_2 \times Z_{2^{n-2}}\)，特别地它不是循环群，但有一个指数为 2 的循环子群（参见 9.5 节推论 20）。

\item 设 \(p\) 是素数，\(V\) 是（用加法书写的）阿贝尔群，具有性质：对所有 \(v \in V\) 有 \(pv = 0\). 若 \(|V| = p^n\)，则 \(V\) 是域 \(\mathbb{F}_p = \mathbb{Z}/p\mathbb{Z}\) 上的 \(n\) 维向量空间。\(V\) 的自同构恰好是 \(V\) 到自身的非奇异线性变换，即
\[
\operatorname{Aut}(V) \cong GL(V) \cong GL_n(\mathbb{F}_p).
\]
特别地，\(\operatorname{Aut}(V)\) 的阶如第 1.4 节所述（参见 10.2 节和 11.1 节的例子）。

\item 对所有 \(n \neq 6\) 有 \(\operatorname{Aut}(S_n) = \operatorname{Inn}(S_n) \cong S_n\)（参见习题 18）。对 \(n = 6\) 有 \(|\operatorname{Aut}(S_6) : \operatorname{Inn}(S_6)| = 2\)（参见下面的习题 19 以及 6.3 节习题 10）。

\item \(\operatorname{Aut}(D_8) \cong D_8\) 且 \(\operatorname{Aut}(Q_8) \cong S_4\)（参见下面的习题 4 和 5 以及 6.3 节习题 9）。
\end{enumerate}
\end{proposition}

命题第 3 部分中描述的群 \(V\) 称为 \(p^n\) 阶\textbf{初等阿贝尔群}（我们将在第 5 章看到，它由 \(p\) 和 \(n\) 在相差同构的意义下唯一确定）。克莱因四元群 \(V_4\) 是 4 阶初等阿贝尔群。这个命题断言
\[
\operatorname{Aut}(V_4) \cong GL_2(\mathbb{F}_2).
\]
由第 1.4 节的习题，后一个群的阶为 6. 但 \(\operatorname{Aut}(V_4)\) 置换 \(V_4\) 的 3 个非单位元，而 \(\operatorname{Aut}(V_4)\) 在 \(V_4 - \{1\}\) 上的这个作用给出 \(\operatorname{Aut}(V_4)\) 到 \(S_3\) 的一个单的置换表示。由阶的考虑，这个同态是满射，故
\[
\operatorname{Aut}(V_4) \cong GL_2(\mathbb{F}_2) \cong S_3.
\]
注意 \(V_4\) 是阿贝尔群，所以 \(\operatorname{Inn}(V_4) = 1\).

对任意素数 \(p\)，\(p^2\) 阶初等阿贝尔群是 \(Z_p \times Z_p\). 它的自同构群 \(GL_2(\mathbb{F}_p)\) 的阶为 \(p(p-1)(p+1)\). 于是推论 \ref{cor:4.9} 蕴含对素数 \(p\)，
\[
\text{若 } |P| = p^2,\quad \text{则 } |\operatorname{Aut}(P)| = p(p-1) \text{ 或 } p(p-1)(p+1),
\]
分别取决于 \(P\) 是循环的或初等阿贝尔的。

\begin{example}
假设 \(G\) 是 \(45 = 3^2 \cdot 5\) 阶群，有一个 \(3^2\) 阶正规子群 \(P\). 我们证明 \(G\) 必是阿贝尔群。

由推论 \ref{cor:4.15}，商 \(G/C_G(P)\) 同构于 \(\operatorname{Aut}(P)\) 的一个子群，而由上一段，\(\operatorname{Aut}(P)\) 的阶为 6 或 48（分别取决于 \(P\) 是循环的或初等阿贝尔的）。另一方面，由于 \(P\) 的阶是素数的平方，\(P\) 是阿贝尔群，故 \(P \leq C_G(P)\). 由此推出 \(|C_G(P)|\) 被 9 整除，这意味着 \(|G/C_G(P)|\) 是 1 或 5. 结合两者得 \(|G/C_G(P)| = 1\)，即 \(C_G(P) = G\) 且 \(P \leq Z(G)\). 由于此时 \(G/Z(G)\) 是循环群，\(G\) 必是阿贝尔群。
\end{example}

\subsection*{习题}
设 \(G\) 是群。

\begin{enumerate}
\item 若 \(\sigma \in \operatorname{Aut}(G)\)，\(\varphi_g\) 是被 \(g\) 共轭，证明 \(\sigma\varphi_g\sigma^{-1} = \varphi_{\sigma(g)}\). 由此推出 \(\operatorname{Inn}(G) \trianglelefteq \operatorname{Aut}(G)\).

\item 证明：若 \(G\) 是 \(pq\) 阶阿贝尔群（\(p, q\) 为不同素数），则 \(G\) 是循环群。〔用柯西定理构造 \(p\) 阶和 \(q\) 阶元素，并考虑它们乘积的阶。〕

\item 证明在 \(D_8\) 的任意自同构下，\(r\) 至多有 2 个可能的像，\(s\) 至多有 4 个可能的像（\(r\) 和 \(s\) 是通常生成元——参见第 1.2 节）。由此推出 \(|\operatorname{Aut}(D_8)| \leq 8\).

\item 用类似前一习题的论证证明 \(|\operatorname{Aut}(Q_8)| \leq 24\).

\item 利用 \(D_8 \trianglelefteq D_{16}\) 这一事实证明 \(\operatorname{Aut}(D_8) \cong D_8\).

\item 证明特征子群是正规的。举出一个正规但非特征子群的例子。

\item 若 \(H\) 是群 \(G\) 中给定阶的唯一子群，证明 \(H\) 在 \(G\) 中特征。

\item 设 \(G\) 是含子群 \(H \leq K\) 的群。

（a）证明：若 \(H\) 在 \(K\) 中特征且 \(K\) 在 \(G\) 中正规，则 \(H\) 在 \(G\) 中正规。

（b）证明：若 \(H\) 在 \(K\) 中特征且 \(K\) 在 \(G\) 中特征，则 \(H\) 在 \(G\) 中特征。用它证明克莱因四元群 \(V_4\) 在 \(S_4\) 中特征。

（c）举一个例子说明：若 \(H\) 在 \(K\) 中正规且 \(K\) 在 \(G\) 中特征，则 \(H\) 不必在 \(G\) 中正规。

\item 设 \(r, s\) 是二面体群 \(D_{2n}\) 的通常生成元。用前两个习题推出 \(\langle r \rangle\) 的每个子群在 \(D_{2n}\) 中都正规。

\item 设 \(G\) 是群，\(A\) 是 \(G\) 的阿贝尔正规子群，记 \(\bar{G} = G/A\). 证明 \(\bar{G}\)（左）通过共轭作用在 \(A\) 上，为 \(\bar{g} \cdot a = gag^{-1}\)，其中 \(g\) 是陪集 \(\bar{g}\) 的任意代表元（特别地，证明这个作用是良定义的）。举一个具体例子说明若 \(A\) 非阿贝尔，则这个作用不是良定义的。

\item 若 \(p\) 是素数，\(P\) 是 \(S_p\) 的 \(p\) 阶子群，证明 \(N_{S_p}(P)/C_{S_p}(P) \cong \operatorname{Aut}(P)\).〔使用第 4.3 节习题 34。〕

\item 设 \(G\) 是 3825 阶群。证明：若 \(H\) 是 \(G\) 的 17 阶正规子群，则 \(H \leq Z(G)\).

\item 设 \(G\) 是 203 阶群。证明：若 \(H\) 是 \(G\) 的 7 阶正规子群，则 \(H \leq Z(G)\). 由此推出此时 \(G\) 是阿贝尔群。

\item 设 \(G\) 是 1575 阶群。证明：若 \(H\) 是 \(G\) 的 9 阶正规子群，则 \(H \leq Z(G)\).

\item 证明下列（乘法）群都是循环的：\((\mathbb{Z}/5\mathbb{Z})^\times\)、\((\mathbb{Z}/9\mathbb{Z})^\times\) 和 \((\mathbb{Z}/18\mathbb{Z})^\times\).

\item 证明 \((\mathbb{Z}/24\mathbb{Z})^\times\) 是 8 阶初等阿贝尔群。（我们将在后面看到 \((\mathbb{Z}/n\mathbb{Z})^\times\) 是初等阿贝尔群当且仅当 \(n \mid 24\).）

\item 设 \(G = \langle x \rangle\) 是 \(n\) 阶循环群。对 \(n = 2, 3, 4, 5, 6\) 具体写出 \(\operatorname{Aut}(G)\) 的元素（由上面命题 \ref{pro:4.16} 我们知道 \(\operatorname{Aut}(G) \cong (\mathbb{Z}/n\mathbb{Z})^\times\)，所以对每个元素 \(a \in (\mathbb{Z}/n\mathbb{Z})^\times\)，具体写出自同构 \(\psi_a\) 对 \(G\) 的元素 \(\{1, x, x^2, \ldots, x^{n-1}\}\) 的作用）。

\item 这个习题表明对 \(n \neq 6\)，\(S_n\) 的每个自同构都是内自同构。固定整数 \(n \geq 2\) 且 \(n \neq 6\).

（a）证明群 \(G\) 的自同构群置换 \(G\) 的共轭类，即对每个 \(\sigma \in \operatorname{Aut}(G)\) 和 \(G\) 的每个共轭类 \(K\)，集合 \(\sigma(K)\) 也是 \(G\) 的共轭类。

（b）设 \(K\) 是 \(S_n\) 中对换的共轭类，\(K'\) 是 \(S_n\) 中任意 2 阶非对换元素的共轭类。证明 \(|K| \neq |K'|\). 由此推出 \(S_n\) 的任何自同构把对换送到对换。〔参见第 3 节习题 33。〕

（c）证明对每个 \(\sigma \in \operatorname{Aut}(S_n)\)，
\[
\sigma((1\,2)) = (a\,b_2), \quad \sigma((1\,3)) = (a\,b_3), \quad \ldots, \quad \sigma((1\,n)) = (a\,b_n)
\]
对某个互不相同的整数 \(a, b_2, b_3, \ldots, b_n \in \{1, 2, \ldots, n\}\).

（d）证明 \((1\,2), (1\,3), \ldots, (1\,n)\) 生成 \(S_n\)，并由此推出 \(S_n\) 的任何自同构由它对这些元素的作用唯一确定。用（c）证明 \(S_n\) 至多有 \(n!\) 个自同构，并得出对 \(n \neq 6\) 有 \(\operatorname{Aut}(S_n) = \operatorname{Inn}(S_n)\).

\item 这个习题表明 \(|\operatorname{Aut}(S_6) : \operatorname{Inn}(S_6)| \leq 2\)（6.3 节习题 10 通过给出一个非内自同构证明了等号成立）。

（a）设 \(K\) 是 \(S_6\) 中对换的共轭类，\(K'\) 是 \(S_6\) 中任意 2 阶非对换元素的共轭类。证明 \(|K| \neq |K'|\)，除非 \(K'\) 是三个不交对换之积的共轭类。由此推出 \(\operatorname{Aut}(S_6)\) 有一个指数至多为 2 的子群把对换送到对换。

（b）证明 \(|\operatorname{Aut}(S_6) : \operatorname{Inn}(S_6)| \leq 2\).〔按前一习题 (c) 和 (d) 的同样步骤证明任何把对换送到对换的自同构都是内自同构。〕

\end{enumerate}

下一个习题引入一个子群 \(J(P)\)，它（像 \(P\) 的中心一样）对任意有限群 \(P\) 都有定义（尽管在大多数应用中 \(P\) 的阶是素数的幂）。这个子群由 J. Thompson 于 1964 年定义，现在在有限群的研究、特别是有限单群的分类中起关键作用。

\begin{enumerate}\setcounter{enumi}{19}

\item 对任意有限群 \(P\)，设 \(d(P)\) 为 \(P\) 的生成元个数的最小值（例如，\(d(P) = 1\) 当且仅当 \(P\) 是非平凡循环群，且 \(d(Q_8) = 2\)）。设 \(m(P)\) 为当 \(A\) 遍历 \(P\) 的所有阿贝尔子群时 \(d(A)\) 的最大值（例如，\(m(Q_8) = 1\)，\(m(D_8) = 2\)）。定义
\[
J(P) = \langle A \mid A \text{ 是 } P \text{ 的阿贝尔子群且 } d(A) = m(P) \rangle.
\]
（\(J(P)\) 称为 \(P\) 的\textbf{汤普森子群}。）

（a）证明 \(J(P)\) 是 \(P\) 的特征子群。

（b）对下列每个群 \(P\)，列出 \(P\) 的所有满足 \(d(A) = m(P)\) 的阿贝尔子群 \(A\)：\(Q_8\)、\(D_8\)、\(D_{16}\) 和 \(QD_{16}\)（其中 \(QD_{16}\) 是第 2.5 节习题 11 中定义的 16 阶拟二面体群）。〔使用第 2.5 节中这些群的子群格。〕

（c）证明 \(J(Q_8) = Q_8\)、\(J(D_8) = D_8\)、\(J(D_{16}) = D_{16}\)，且 \(J(QD_{16})\) 是 \(QD_{16}\) 中的 8 阶二面体子群。

（d）证明：若 \(Q \leq P\) 且 \(J(P)\) 是 \(Q\) 的子群，则 \(J(P) = J(Q)\). 由此推出：若 \(P\) 是有限群 \(G\) 的子群（不必正规），且 \(J(P)\) 包含在 \(P\) 的某个满足 \(Q \trianglelefteq G\) 的子群 \(Q\) 中，则 \(J(P) \trianglelefteq G\).
\end{enumerate}

\section{西罗定理}

本节我们证明拉格朗日定理的一个部分逆定理，并推出大量推论，其中一些将在下一章导向分类定理。

\begin{definition}
设 \(G\) 是群，\(p\) 是素数。

（1）阶为 \(p^\alpha\)（\(\alpha \geq 1\)）的群称为 \textbf{\(p\)-群}。\(G\) 的 \(p\)-子群是指本身是 \(p\)-群的子群。

（2）若 \(G\) 是 \(p^\alpha m\) 阶群（\(p \nmid m\)），则阶为 \(p^\alpha\) 的子群称为 \(G\) 的\textbf{西罗 \(p\)-子群}。

（3）\(G\) 的西罗 \(p\)-子群的集合记作 \(\operatorname{Syl}_p(G)\)，\(G\) 的西罗 \(p\)-子群的个数记作 \(n_p(G)\)（或当 \(G\) 由上下文清楚时简记为 \(n_p\)）。
\end{definition}

\begin{theorem}[西罗定理]\label{thm:4.18}
设 \(G\) 是 \(p^\alpha m\) 阶群，其中 \(p\) 是素数且不整除 \(m\).

\begin{enumerate}[label=(\arabic*)]
\item \(G\) 的西罗 \(p\)-子群存在，即 \(\operatorname{Syl}_p(G) \neq \varnothing\).

\item 若 \(P\) 是 \(G\) 的西罗 \(p\)-子群，\(Q\) 是 \(G\) 的任意 \(p\)-子群，则存在 \(g \in G\) 使得 \(Q \leq gPg^{-1}\)，即 \(Q\) 包含在 \(P\) 的某个共轭中。特别地，\(G\) 的任意两个西罗 \(p\)-子群在 \(G\) 中共轭。

\item \(G\) 的西罗 \(p\)-子群的个数 \(n_p\) 形如 \(1 + kp\)，即 \(n_p \equiv 1 \pmod{p}\). 此外，对任意西罗 \(p\)-子群 \(P\)，\(n_p\) 是 \(N_G(P)\) 在 \(G\) 中的指数，因而整除 \(m\).
\end{enumerate}
\end{theorem}

我们首先证明下面的引理：

\begin{lemma}\label{lem:4.19}
设 \(P \in \operatorname{Syl}_p(G)\). 若 \(Q\) 是 \(G\) 的任意 \(p\)-子群，则 \(Q \cap N_G(P) = Q \cap P\).
\end{lemma}

\begin{proof}
设 \(H = N_G(P) \cap Q\). 由于 \(P \leq N_G(P)\)，显然 \(P \cap Q \leq H\)，所以我们要证明反向包含。由于按定义 \(H \leq Q\)，这等价于证明 \(H \leq P\). 我们通过证明 \(PH\) 是 \(G\) 的一个同时包含 \(P\) 和 \(H\) 的 \(p\)-子群来做到这一点；但 \(P\) 是 \(G\) 的阶尽可能大的 \(p\)-子群，故必有 \(PH = P\)，即 \(H \leq P\).

由于 \(H \leq N_G(P)\)，由第 3.2 节推论 \ref{cor:3.15}，\(PH\) 是子群。由同节的命题 \ref{pro:3.13}，
\[
|PH| = \frac{|P||H|}{|P \cap H|}.
\]
上述商中所有数都是 \(p\) 的幂，故 \(PH\) 是 \(p\)-群。此外，\(P\) 是 \(PH\) 的子群，故 \(PH\) 的阶被 \(p^\alpha\)（整除 \(|G|\) 的 \(p\) 的最大幂）整除。这两个事实迫使 \(|PH| = p^\alpha = |P|\). 这又蕴含 \(P = PH\) 且 \(H \leq P\). 引理得证。
\end{proof}

\begin{proof}[西罗定理的证明]
（1）对 \(|G|\) 作归纳。若 \(|G| = 1\)，无需证明。归纳假设所有阶小于 \(|G|\) 的群都存在西罗 \(p\)-子群。

若 \(p\) 整除 \(|Z(G)|\)，则由阿贝尔群的柯西定理（第 3.4 节命题 \ref{pro:3.21}）\(Z(G)\) 有 \(p\) 阶子群 \(N\). 令 \(\bar{G} = G/N\)，则 \(|\bar{G}| = p^{\alpha-1}m\). 由归纳假设，\(\bar{G}\) 有 \(p^{\alpha-1}\) 阶子群 \(\bar{P}\). 设 \(P\) 是 \(G\) 中包含 \(N\) 且 \(P/N = \bar{P}\) 的子群，则 \(|P| = |P/N| \cdot |N| = p^\alpha\)，且 \(P\) 是 \(G\) 的西罗 \(p\)-子群。于是我们化归到 \(p\) 不整除 \(|Z(G)|\) 的情形。

设 \(g_1, g_2, \ldots, g_r\) 是 \(G\) 的不同非中心共轭类的代表元。\(G\) 的类方程是
\[
|G| = |Z(G)| + \sum_{i=1}^{r} |G : C_G(g_i)|.
\]
若对所有 \(i\) 都有 \(p \mid |G : C_G(g_i)|\)，则由于 \(p \mid |G|\)，也会有 \(p \mid |Z(G)|\)，矛盾。因此对某个 \(i\)，\(p\) 不整除 \(|G : C_G(g_i)|\). 对这个 \(i\)，令 \(H = C_G(g_i)\)，则
\[
|H| = p^\alpha k, \quad \text{其中 } p \nmid k.
\]
由于 \(g_i \notin Z(G)\)，\(|H| < |G|\). 由归纳假设，\(H\) 有西罗 \(p\)-子群 \(P\)，它当然也是 \(G\) 的子群。由于 \(|P| = p^\alpha\)，\(P\) 是 \(G\) 的西罗 \(p\)-子群。这完成了归纳并证明了（1）。

在证明（2）和（3）之前我们作一些计算。由（1）存在 \(G\) 的西罗 \(p\)-子群 \(P\). 设
\[
S = \{P_1, P_2, \ldots, P_r\}
\]
是 \(P\) 的所有共轭的集合（即 \(S = \{gPg^{-1} \mid g \in G\}\)），设 \(Q\) 是 \(G\) 的任意 \(p\)-子群。由 \(S\) 的定义，\(G\)（从而 \(Q\)）通过共轭作用在 \(S\) 上。把这个作用下的 \(S\) 写成轨道的并：
\[
S = O_1 \cup O_2 \cup \cdots \cup O_s,
\]
其中 \(r = |O_1| + \cdots + |O_s|\). 请记住 \(r\) 不依赖 \(Q\)，但 \(Q\)-轨道的个数 \(s\) 依赖 \(Q\)（注意按定义 \(G\) 在 \(S\) 上只有一个轨道，但 \(G\) 的子群 \(Q\) 可能有不止一个轨道）。必要时重新给 \(S\) 的元素编号，使 \(S\) 的前 \(s\) 个元素是 \(Q\)-轨道的代表元：\(P_i \in O_i\)，\(1 \leq i \leq s\). 由命题 \ref{pro:4.2} 可得 \(|O_i| = |Q : N_Q(P_i)|\). 按定义 \(N_Q(P_i) = N_G(P_i) \cap Q\)，由引理 \ref{lem:4.19}，\(N_G(P_i) \cap Q = P_i \cap Q\). 结合这两个事实得
\[
|O_i| = |Q : P_i \cap Q|, \quad 1 \leq i \leq s. \tag{4.1}
\]
现在我们可以证明 \(r \equiv 1 \pmod{p}\) 了。由于 \(Q\) 是任意的，我们可以在上面取 \(Q = P_1\)，于是 (4.1) 给出
\[
|O_1| = |P_1 : P_1 \cap P_1| = 1.
\]
现在对所有 \(i > 1\) 有 \(P_1 \neq P_i\)，故 \(P_1 \cap P_i < P_1\). 由 (4.1)，
\[
|O_i| = |P_1 : P_1 \cap P_i| > 1, \quad 2 \leq i \leq s.
\]
由于 \(P_1\) 是 \(p\)-群，\(|P_1 : P_1 \cap P_i|\) 必是 \(p\) 的幂，故
\[
p \mid |O_i|, \quad 2 \leq i \leq s.
\]
于是
\[
r = |O_1| + (|O_2| + \cdots + |O_s|) \equiv 1 \pmod{p}.
\]

现在我们证明（2）和（3）。设 \(Q\) 是 \(G\) 的任意 \(p\)-子群。假设 \(Q\) 不含于任何 \(P_i\) 中，即对所有 \(i \in \{1, 2, \ldots, r\}\) 有 \(Q \not\leq P_i\)（即对任意 \(g \in G\) 有 \(Q \not\leq gPg^{-1}\)）。此时对每个 \(i\)，由 (4.1) 有 \(|O_i| = |Q : Q \cap P_i| > 1\). 于是对所有 \(i\) 有 \(p \mid |O_i|\)，故 \(p\) 整除 \(|O_1| + \cdots + |O_s| = r\). 这与 \(r \equiv 1 \pmod{p}\) 矛盾（记住 \(r\) 不依赖 \(Q\) 的选取）。这个矛盾证明了 \(Q \leq gPg^{-1}\) 对某个 \(g \in G\) 成立。

为看出 \(G\) 的所有西罗 \(p\)-子群共轭，设 \(Q\) 是 \(G\) 的任意西罗 \(p\)-子群。由前面的论证，对某个 \(g \in G\) 有 \(Q \leq gPg^{-1}\). 由于 \(|gPg^{-1}| = |Q| = p^\alpha\)，必有 \(gPg^{-1} = Q\). 这证明了定理的（2）。特别地，\(S = \operatorname{Syl}_p(G)\)，因为 \(G\) 的每个西罗 \(p\)-子群都与 \(P\) 共轭，故 \(n_p = r \equiv 1 \pmod{p}\)，这是（3）的第一部分。

最后，由于所有西罗 \(p\)-子群共轭，命题 \ref{pro:4.6} 表明对任意 \(P \in \operatorname{Syl}_p(G)\) 有
\[
n_p = |G : N_G(P)|,
\]
完成西罗定理的证明。
\end{proof}

注意西罗定理的共轭部分连同推论 \ref{cor:4.14} 表明群（对同一素数 \(p\)）的任意两个西罗 \(p\)-子群同构。

\begin{corollary}\label{cor:4.20}
设 \(P\) 是 \(G\) 的西罗 \(p\)-子群。则下列命题等价：

\begin{enumerate}[label=(\arabic*)]
\item \(P\) 是 \(G\) 的唯一的西罗 \(p\)-子群，即 \(n_p = 1\)；
\item \(P\) 在 \(G\) 中正规；
\item \(P\) 在 \(G\) 中特征；
\item 由 \(p\)-幂阶元素生成的所有子群都是 \(p\)-群，即若 \(X\) 是 \(G\) 的任意子集，使得对所有 \(x \in X\) 有 \(|x|\) 是 \(p\) 的幂，则 \(\langle X \rangle\) 是 \(p\)-群。
\end{enumerate}
\end{corollary}

\begin{proof}
若（1）成立，则对所有 \(g \in G\) 有 \(gPg^{-1} = P\)，因为 \(gPg^{-1} \in \operatorname{Syl}_p(G)\)，即 \(P\) 在 \(G\) 中正规。故（1）蕴含（2）。反之，若 \(P \trianglelefteq G\) 且 \(Q \in \operatorname{Syl}_p(G)\)，则由西罗定理存在 \(g \in G\) 使 \(Q = gPg^{-1} = P\). 故 \(\operatorname{Syl}_p(G) = \{P\}\)，（2）蕴含（1）。

由于特征子群正规，（3）蕴含（2）。反之，若 \(P \trianglelefteq G\)，我们刚刚证明 \(P\) 是 \(G\) 中唯一的 \(p^\alpha\) 阶子群，故 \(P\ \mathrm{char}\ G\). 于是（2）与（3）等价。

最后，假设（1）成立，并设 \(X\) 是 \(G\) 的子集，使得对所有 \(x \in X\) 有 \(|x|\) 是 \(p\) 的幂。由西罗定理的共轭部分，对每个 \(x \in X\) 存在某个 \(g \in G\) 使 \(x \in gPg^{-1} = P\). 于是 \(X \subseteq P\)，从而 \(\langle X \rangle \leq P\)，\(\langle X \rangle\) 是 \(p\)-群。反之，若（4）成立，设 \(X\) 是 \(G\) 的所有西罗 \(p\)-子群的并。若 \(P\) 是任意西罗 \(p\)-子群，则 \(P\) 是 \(p\)-群 \(\langle X \rangle\) 的子群。由于 \(P\) 是 \(G\) 的阶最大的 \(p\)-子群，必有 \(P = \langle X \rangle\)，故（1）成立。
\end{proof}

\begin{example}
设 \(G\) 是有限群，\(p\) 是素数。

（1）若 \(p\) 不整除 \(G\) 的阶，则 \(G\) 的西罗 \(p\)-子群是平凡群（西罗定理的所有部分都平凡地成立）。若 \(|G| = p^\alpha\)，则 \(G\) 是 \(G\) 的唯一的西罗 \(p\)-子群。

（2）有限阿贝尔群对每个素数 \(p\) 都有唯一的西罗 \(p\)-子群。这个子群由所有阶为 \(p\) 的幂的元素组成，有时称为该阿贝尔群的 \textbf{\(p\)-准素分量}。

（3）\(S_3\) 有三个西罗 2-子群：\(\langle (1\,2) \rangle\)、\(\langle (2\,3) \rangle\) 和 \(\langle (1\,3) \rangle\). 它有唯一的（因此正规的）西罗 3-子群：\(\langle (1\,2\,3) \rangle = A_3\). 注意 \(3 \equiv 1 \pmod{2}\).

（4）\(A_4\) 有唯一的西罗 2-子群：\(\langle (1\,2)(3\,4), (1\,3)(2\,4) \rangle \cong V_4\). 它有四个西罗 3-子群：\(\langle (1\,2\,3) \rangle\)、\(\langle (1\,2\,4) \rangle\)、\(\langle (1\,3\,4) \rangle\) 和 \(\langle (2\,3\,4) \rangle\). 注意 \(4 \equiv 1 \pmod{3}\).

（5）\(S_4\) 有 \(n_2 = 3\) 和 \(n_3 = 4\). 由于 \(S_4\) 含有一个同构于 \(D_8\) 的子群，\(S_4\) 的每个西罗 2-子群都同构于 \(D_8\).
\end{example}

\subsection*{西罗定理的应用}

我们现在给出西罗定理的一些应用。大多数例子用西罗定理证明某个特定阶的群不是单群。在讨论了从较小群构造较大群的方法（例如半直积的构造）之后，我们将能利用这些结果来分类某些特定阶 \(n\) 的群（正如我们已经对 \(n = 15\) 所做的那样）。

由于西罗定理保证了有限群的 \(p\)-子群的存在，更细致地研究素数幂阶群是值得的。这将在第 6 章进行，那里将讨论西罗定理的更多应用。

对小阶群，仅西罗定理的同余条件往往就足以迫使非平凡正规子群存在。西罗定理任何数值应用的第一步是把群的阶分解成素数幂。群阶的最大素因子往往给出 \(n_p\) 最少的可能取值（例如，关于 \(n_2\) 的同余条件不给出任何限制），这限制了群 \(G\) 的结构。在下面的例子中我们将看到仅西罗定理不能迫使正规子群存在的情形，然而某个附加论证（常常涉及对若干个不同素数 \(p\) 的 \(p\) 阶元素的研究）证明了正规西罗子群的存在。

\begin{example}[\(pq\) 阶群，\(p, q\) 是素数且 \(p < q\)]
假设 \(|G| = pq\)，其中 \(p, q\) 是素数且 \(p < q\). 设 \(P \in \operatorname{Syl}_p(G)\)，\(Q \in \operatorname{Syl}_q(G)\). 我们证明 \(Q\) 在 \(G\) 中正规，且若 \(P\) 也在 \(G\) 中正规，则 \(G\) 是循环群。

现在三个条件：\(n_q = 1 + kq\)（\(k \geq 0\)）、\(n_q\) 整除 \(p\)、\(p < q\)，共同迫使 \(k = 0\). 由于 \(n_q = 1\)，\(Q \trianglelefteq G\).

由于 \(n_p\) 整除素数 \(q\)，只可能 \(n_p = 1\) 或 \(q\). 特别地，若 \(p \nmid q - 1\)（即若 \(q \not\equiv 1 \pmod{p}\)），则 \(n_p\) 不能等于 \(q\)，故 \(P \trianglelefteq G\).

设 \(P = \langle x \rangle\)、\(Q = \langle y \rangle\). 若 \(P \trianglelefteq G\)，则由于 \(G/C_G(P)\) 同构于 \(\operatorname{Aut}(Z_p)\) 的一个子群，而后者的阶为 \(p - 1\)，拉格朗日定理连同 \(p\) 和 \(q\) 都不整除 \(p - 1\) 这一观察蕴含 \(G = C_G(P)\). 此时 \(x \in P \leq Z(G)\)，所以 \(x\) 与 \(y\) 交换。（或者，这立即由第 3.1 节习题 42 推出。）这意味着 \(|xy| = pq\)（参见第 2.3 节的习题），故此时 \(G\) 是循环群：\(G \cong Z_{pq}\).

若 \(p \mid q - 1\)，我们将在第 5 章看到存在唯一的 \(pq\) 阶非阿贝尔群（其中必有 \(n_p = q\)）。我们现在就可以证明这个群的存在性。设 \(Q\) 是 \(q\) 次对称群 \(S_q\) 的西罗 \(q\)-子群。由第 3 节习题 34，\(|N_{S_q}(Q)| = q(q-1)\). 由假设 \(p \mid q-1\)，故由柯西定理 \(N_{S_q}(Q)\) 有 \(p\) 阶子群 \(P\). 由第 3.2 节推论 \ref{cor:3.15}，\(PQ\) 是 \(pq\) 阶群。由于 \(C_{S_q}(Q) = Q\)（第 3 节例 2），\(PQ\) 是非阿贝尔群。\(PQ\) 唯一性证明的本质要素是关于 \(\operatorname{Aut}(Z_q)\) 循环性的命题 \ref{pro:4.17}.
\end{example}

\begin{example}[30 阶群]
设 \(G\) 是 30 阶群。我们证明 \(G\) 有同构于 \(Z_{15}\) 的正规子群。我们将利用这个信息在下一章分类 30 阶群。注意任何 15 阶子群必正规（因其指数为 2）且循环（由前面的结果），所以只需证明存在 15 阶子群。最快的办法是引用第 2 节习题 13. 我们给出另一个论证，说明如何把西罗定理与素数阶元素的计数结合起来产生正规子群。

设 \(P \in \operatorname{Syl}_5(G)\)，\(Q \in \operatorname{Syl}_3(G)\). 若 \(P\) 或 \(Q\) 中有一个在 \(G\) 中正规，则由第 3 章推论 \ref{cor:3.15}，\(PQ\) 是 15 阶群。还要注意若 \(P\) 或 \(Q\) 中有一个正规，则 \(P\) 和 \(Q\) 都是 \(PQ\) 的特征子群，且由于 \(PQ \trianglelefteq G\)，\(P\) 和 \(Q\) 都在 \(G\) 中正规（第 4 节习题 8(a)）。因此假设两个西罗子群都不正规。由西罗定理第 3 部分，唯一的可能是 \(n_5 = 6\) 和 \(n_3 = 10\).

每个 5 阶元素都落在一个西罗 5-子群中，每个西罗 5-子群含 4 个非单位元，且由拉格朗日定理，不同的西罗 5-子群交为单位元。因此 \(G\) 中 5 阶元素的个数等于一个西罗 5-子群中非单位元的个数乘以西罗 5-子群的个数，这将是 \(4 \cdot 6 = 24\) 个 5 阶元素。由类似的推理，3 阶元素的个数将是 \(2 \cdot 10 = 20\). 这很荒谬，因为 30 阶群不可能含有 \(24 + 20 = 44\) 个不同元素。\(P\) 或 \(Q\) 中必有一个（从而两者都）在 \(G\) 中正规。
\end{example}

这种计数技巧经常有用（另见 6.2 节），并且当西罗 \(p\)-子群的阶为 \(p\) 时（如本例）特别有效，因为此时两个不同西罗 \(p\)-子群的交必为单位元。若西罗 \(p\)-子群的阶为 \(p^\alpha\)（\(\alpha \geq 2\)），计数元素时需要更加小心，因为此时不同的西罗 \(p\)-子群可能有更多公共元素，即交可能非平凡。

\begin{example}[12 阶群]
设 \(G\) 是 12 阶群。我们证明或者 \(G\) 有正规的西罗 3-子群，或者 \(G \cong A_4\)（在后一情形 \(G\) 有正规的西罗 2-子群）。我们将利用这个信息在下一章分类 12 阶群。

假设 \(n_3 \neq 1\)，设 \(P \in \operatorname{Syl}_3(G)\). 由于 \(n_3 \mid 4\) 且 \(n_3 \equiv 1 \pmod{3}\)，可知 \(n_3 = 4\). 由于不同西罗 3-子群交为单位元且每个含两个 3 阶元素，\(G\) 含 \(2 \cdot 4 = 8\) 个 3 阶元素。由于 \(|G : N_G(P)| = n_3 = 4\)，\(N_G(P) = P\). 现在 \(G\) 通过共轭作用在其四个西罗 3-子群上，所以这个作用提供一个置换表示
\[
\varphi : G \to S_4
\]
（注意我们也可以作用在 \(P\) 的左陪集上并用定理 \ref{thm:4.3}）。这个作用的核 \(K\) 是 \(G\) 中正规化 \(G\) 的所有西罗 3-子群的子群。特别地，\(K \leq N_G(P) = P\). 由于由假设 \(P\) 在 \(G\) 中不正规，\(K = 1\)，即 \(\varphi\) 是单射且
\[
|G| = |\varphi(G)| = 12.
\]
由于 \(G\) 含 8 个 3 阶元素，而 \(S_4\) 中恰好有 8 个 3 阶元素（都落在 \(A_4\) 中），可知 \(\varphi(G)\) 与 \(A_4\) 的交是阶至少为 8 的子群。由于两个群的阶都是 12，可知 \(\varphi(G) = A_4\)，故 \(G \cong A_4\).

注意 \(A_4\) 确实有 4 个西罗 3-子群（见推论 \ref{cor:4.20} 后的例 4），所以这样的群 \(G\) 确实存在。此外，设 \(V\) 是 \(A_4\) 的西罗 2-子群。由于 \(|V| = 4\)，它包含 \(A_4\) 的所有其余元素。特别地，不可能有另一个西罗 2-子群。于是 \(n_2(A_4) = 1\)，即 \(V \trianglelefteq A_4\)（也可以直接看出，因为 \(V\) 是单位元连同 \(A_4\) 中三个两两不交对换之积的元素，即 \(V\) 是共轭类的并）。
\end{example}

\begin{example}[\(p^2 q\) 阶群，\(p, q\) 是不同素数]
设 \(G\) 是 \(p^2 q\) 阶群。我们证明 \(G\) 有正规的西罗子群（对 \(p\) 或 \(q\) 之一）。我们将利用这个信息在下一章分类这种阶的一些群（参见 5.5 节习题 8 至 12）。设 \(P \in \operatorname{Syl}_p(G)\)，\(Q \in \operatorname{Syl}_q(G)\).

首先考虑 \(p > q\) 的情形。由于 \(n_p \mid q\) 且 \(n_p = 1 + kp\)，必有 \(n_p = 1\). 故 \(P \trianglelefteq G\).

现在考虑 \(p < q\) 的情形。若 \(n_q = 1\)，则 \(Q\) 在 \(G\) 中正规。因此假设 \(n_q > 1\)，即对某个 \(t > 0\) 有 \(n_q = 1 + tq\). 现在 \(n_q\) 整除 \(p^2\)，故 \(n_q = p\) 或 \(p^2\). 由于 \(q > p\)，不可能有 \(n_q = p\)，故 \(n_q = p^2\). 于是
\[
tq = p^2 - 1 = (p-1)(p+1).
\]
由于 \(q\) 是素数，或者 \(q \mid p-1\)，或者 \(q \mid p+1\). 前者不可能，因为 \(q > p\)，故后者成立。由于 \(q > p\) 但 \(q \mid p+1\)，必有 \(q = p + 1\). 这迫使 \(p = 2\)、\(q = 3\) 且 \(|G| = 12\). 结果随即由前一个例子得出。
\end{example}

\subsection*{60 阶群}

我们说明如何用西罗定理揭示给定阶群的结构，即使该阶的某些群可能是单群。注意从一个素数切换到另一个素数的技巧，以及利用阶小于 60 的群的结果来研究 60 阶群的归纳过程。

\begin{proposition}\label{pro:4.21}
若 \(|G| = 60\) 且 \(G\) 有不止一个西罗 5-子群，则 \(G\) 是单群。
\end{proposition}

\begin{proof}
反设 \(|G| = 60\) 且 \(n_5 > 1\)，但存在正规子群 \(H\) 满足 \(H \neq 1, G\). 由西罗定理，\(n_5\) 的唯一可能是 6. 设 \(P \in \operatorname{Syl}_5(G)\)，则 \(|N_G(P)| = 10\)，因为其指数为 \(n_5\).

若 \(5 \mid |H|\)，则 \(H\) 含 \(G\) 的一个西罗 5-子群，且由于 \(H\) 正规，它包含这个子群的所有 6 个共轭。特别地，\(|H| \geq 1 + 6 \cdot 4 = 25\)，唯一可能是 \(|H| = 30\). 这导致矛盾，因为前面一个例子证明了任何 30 阶群都有正规的（因此唯一的）西罗 5-子群。这个论证表明 \(5\) 不整除 \(H\) 对 \(G\) 的任何真正规子群都成立。

若 \(|H| = 6\) 或 12，则 \(H\) 有正规的、因而特征的西罗子群，它在 \(G\) 中也正规。必要时把 \(H\) 替换为这个子群，我们可以假设 \(|H| = 2, 3\) 或 4. 令 \(\bar{G} = G/H\)，则 \(|\bar{G}| = 30, 20\) 或 15. 在每种情形，由前面的结果 \(\bar{G}\) 有 5 阶正规子群 \(\bar{P}\). 设 \(H_1\) 是 \(\bar{P}\) 在 \(G\) 中的完全原像，则 \(H_1 \trianglelefteq G\)、\(H_1 \neq 1\) 且 \(5 \mid |H_1|\). 这与前一段矛盾，从而完成证明。
\end{proof}

\begin{corollary}\label{cor:4.22}
\(A_5\) 是单群。
\end{corollary}

\begin{proof}
子群 \(\langle (1\,2\,3\,4\,5) \rangle\) 和 \(\langle (1\,3\,2\,4\,5) \rangle\) 是 \(A_5\) 的两个不同西罗 5-子群，故结果立即由命题 \ref{pro:4.21} 得出。
\end{proof}

下一个命题表明 60 阶单群只有一个。

\begin{proposition}\label{pro:4.23}
若 \(G\) 是 60 阶单群，则 \(G \cong A_5\).
\end{proposition}

\begin{proof}
设 \(G\) 是 60 阶单群，则 \(n_2 = 3, 5\) 或 15. 设 \(P \in \operatorname{Syl}_2(G)\)，\(N = N_G(P)\)，故 \(|G : N| = n_2\).

首先注意 \(G\) 没有指数小于 5 的真子群，理由如下：若 \(H\) 是 \(G\) 的指数为 4、3 或 2 的子群，则由定理 \ref{thm:4.3}，\(G\) 有包含在 \(H\) 中的正规子群 \(K\)，且 \(G/K\) 同构于 \(S_4\)、\(S_3\) 或 \(S_2\) 的一个子群。由于 \(K \neq G\)，单性迫使 \(K = 1\). 这不可能，因为 60（\(= |G|\)）不整除 4!. 这个论证特别表明 \(n_2 \neq 3\).

若 \(n_2 = 5\)，则 \(N\) 在 \(G\) 中指数为 5，故 \(G\) 通过左乘作用在 \(N\) 的左陪集集合上给出 \(G\) 到 \(S_5\) 的置换表示。由于（如上）这个表示的核是真正规子群而 \(G\) 是单群，核是 1，且 \(G\) 同构于 \(S_5\) 的一个子群。把 \(G\) 与这个同构复制等同起来，于是可设 \(G \leq S_5\). 若 \(G\) 不含于 \(A_5\)，则 \(S_5 = GA_5\)，由第二同构定理，\(A_5 \cap G\) 在 \(G\) 中指数为 2. 由于 \(G\) 没有指数为 2 的（正规）子群，这是矛盾的。这个论证证明了 \(G \cong A_5\). 由于 \(|G| = |A_5|\)，\(G\) 在 \(S_5\) 中的同构复制与 \(A_5\) 重合，正如所要。

最后，假设 \(n_2 = 15\). 若对 \(G\) 的每一对不同的西罗 2-子群 \(P, Q\) 都有 \(P \cap Q = 1\)，则 \(G\) 的西罗 2-子群中非单位元的个数为 \((4-1) \cdot 15 = 45\). 但 \(n_5 = 6\)，故 \(G\) 中 5 阶元素的个数为 \((5-1) \cdot 6 = 24\)，合计 69 个元素。这个矛盾证明了存在不同的西罗 2-子群 \(P, Q\) 使 \(|P \cap Q| = 2\). 设 \(M = N_G(P \cap Q)\). 由于 \(P\) 和 \(Q\) 是阿贝尔群（作为 4 阶群），\(P\) 和 \(Q\) 是 \(M\) 的子群，且由于 \(G\) 是单群，\(M \neq G\). 于是 \(4\) 整除 \(|M|\) 且 \(|M| > 4\)（否则 \(P = M = Q\)）。唯一可能是 \(|M| = 12\)，即 \(M\) 在 \(G\) 中指数为 5（回忆 \(M\) 不可能有指数 3 或 1）。但现在把前一段的论证应用于 \(M\)（代替 \(N\)）得到 \(G \cong A_5\). 这在本情形导致矛盾，因为 \(n_2(A_5) = 5\)（参见习题）。证明完成。
\end{proof}

\subsection*{习题}
设 \(G\) 是有限群，\(p\) 是素数。

\begin{enumerate}
\item 证明：若 \(P \in \operatorname{Syl}_p(G)\) 且 \(H\) 是 \(G\) 的包含 \(P\) 的子群，则 \(P \in \operatorname{Syl}_p(H)\). 举一个例子说明一般地，\(G\) 的子群的西罗 \(p\)-子群不必是 \(G\) 的西罗 \(p\)-子群。

\item 证明：若 \(H\) 是 \(G\) 的子群且 \(Q \in \operatorname{Syl}_p(H)\)，则对所有 \(g \in G\) 有 \(gQg^{-1} \in \operatorname{Syl}_p(gHg^{-1})\).

\item 用西罗定理证明柯西定理。（注意我们在西罗定理的证明中只对阿贝尔群使用了柯西定理——命题 \ref{pro:3.21}——所以这条推理路线不是循环论证。）

\item 列出 \(D_{12}\) 和 \(S_3 \times S_3\) 的所有西罗 2-子群和西罗 3-子群。

\item 证明对每个奇素数 \(p\)，\(D_{2n}\) 的西罗 \(p\)-子群都是循环且正规的。

\item 列出 \(A_4\) 的所有西罗 3-子群和 \(S_4\) 的所有西罗 3-子群。

\item 列出 \(S_4\) 的所有西罗 2-子群，并求 \(S_4\) 中把这些子群之一共轭到其他每个子群的元素。

\item 列出 \(S_5\) 的两个不同西罗 2-子群，以及 \(S_5\) 中把其中一个共轭为另一个的元素。

\item 列出 \(SL_2(\mathbb{F}_3)\) 的所有西罗 3-子群（参见第 2.1 节习题 9）。

\item 证明由 \(\begin{pmatrix} 0 & -1 \\ 1 & 0 \end{pmatrix}\) 和 \(\begin{pmatrix} 1 & 1 \\ 1 & -1 \end{pmatrix}\) 生成的 \(SL_2(\mathbb{F}_3)\) 的子群是 \(SL_2(\mathbb{F}_3)\) 的唯一的西罗 2-子群（参见第 2.4 节习题 10）。

\item 证明 \(SL_2(\mathbb{F}_3)\) 的中心是由 \(\pm I\) 组成的 2 阶群，其中 \(I\) 是单位矩阵。证明 \(SL_2(\mathbb{F}_3)/Z(SL_2(\mathbb{F}_3)) \cong A_4\).〔使用关于 12 阶群的事实。〕

\item 设 \(|D_{2n}| = 2n = 2^\alpha k\)，其中 \(k\) 是奇数。证明 \(D_{2n}\) 的西罗 2-子群的个数为 \(k\).〔证明若 \(P \in \operatorname{Syl}_2(D_{2n})\)，则 \(N_{D_{2n}}(P) = P\).〕

\item 证明 56 阶群有正规的西罗 \(p\)-子群，对某个整除其阶的素数 \(p\).

\item 证明 312 阶群有正规的西罗 \(p\)-子群，对某个整除其阶的素数 \(p\).

\item 证明 351 阶群有正规的西罗 \(p\)-子群，对某个整除其阶的素数 \(p\).

\item 设 \(|G| = pqr\)，其中 \(p, q, r\) 是素数且 \(p < q < r\). 证明 \(G\) 有正规的西罗子群，对 \(p, q, r\) 之一。

\item 证明若 \(|G| = 105\)，则 \(G\) 有正规的西罗 5-子群和正规的西罗 7-子群。

\item 证明 200 阶群有正规的西罗 5-子群。

\item 证明若 \(|G| = 6545\)，则 \(G\) 不是单群。

\item 证明若 \(|G| = 1365\)，则 \(G\) 不是单群。

\item 证明若 \(|G| = 2907\)，则 \(G\) 不是单群。

\item 证明若 \(|G| = 132\)，则 \(G\) 不是单群。

\item 证明若 \(|G| = 462\)，则 \(G\) 不是单群。

\item 证明若 \(G\) 是 231 阶群，则 \(Z(G)\) 包含 \(G\) 的一个西罗 11-子群，且 \(G\) 的一个西罗 7-子群正规。

\item 证明若 \(G\) 是 385 阶群，则 \(Z(G)\) 包含 \(G\) 的一个西罗 7-子群，且 \(G\) 的一个西罗 11-子群正规。

\item 设 \(G\) 是 105 阶群。证明：若 \(G\) 的一个西罗 3-子群正规，则 \(G\) 是阿贝尔群。

\item 设 \(G\) 是 315 阶群，有一个正规的西罗 3-子群。证明 \(Z(G)\) 包含 \(G\) 的一个西罗 3-子群，并由此推出 \(G\) 是阿贝尔群。

\item 设 \(G\) 是 1575 阶群。证明：若 \(G\) 的一个西罗 3-子群正规，则 \(G\) 的一个西罗 5-子群和一个西罗 7-子群也正规。在此情形证明 \(G\) 是阿贝尔群。

\item 若 \(G\) 是阶不超过 100 的非阿贝尔单群，证明 \(G \cong A_5\).〔排除除 60 以外的所有阶。〕

\item 168 阶单群中必有多少个 7 阶元素？

\item 对 \(p = 2, 3, 5\) 求 \(n_p(A_5)\) 和 \(n_p(S_5)\).〔注意 \(A_4 \leq A_5\).〕

\item 设 \(P\) 是 \(H\) 的西罗 \(p\)-子群，\(H\) 是 \(K\) 的子群。若 \(P \trianglelefteq H\) 且 \(H \trianglelefteq K\)，证明 \(P\) 在 \(K\) 中正规。由此推出：若 \(P \in \operatorname{Syl}_p(G)\) 且 \(H = N_G(P)\)，则 \(N_G(H) = H\)（即：西罗 \(p\)-子群的正规化子是自正规化的）。

\item 设 \(P\) 是 \(G\) 的正规西罗 \(p\)-子群，\(H\) 是 \(G\) 的任意子群。证明 \(P \cap H\) 是 \(H\) 的唯一的西罗 \(p\)-子群。

\item 设 \(P \in \operatorname{Syl}_p(G)\) 且假设 \(N \trianglelefteq G\). 用西罗定理的共轭部分证明 \(P \cap N\) 是 \(N\) 的西罗 \(p\)-子群。由此推出 \(PN/N\) 是 \(G/N\) 的西罗 \(p\)-子群（注意这也可以用第二同构定理来做——参见第 3.3 节习题 9）。

\item 设 \(P \in \operatorname{Syl}_p(G)\)，\(H \leq G\). 证明对某个 \(g \in G\)，\(gPg^{-1} \cap H\) 是 \(H\) 的西罗 \(p\)-子群。举一个具体例子说明对任意 \(h \in H\)，\(hPh^{-1} \cap H\) 不必是 \(H\) 的西罗 \(p\)-子群（特别地，我们不能总是在第一部分中取 \(g = 1\)，而在 \(H\) 在 \(G\) 中正规时是可以的）。

\item 证明：若 \(N\) 是 \(G\) 的正规子群，则 \(n_p(G/N) \leq n_p(G)\).

\item 设 \(R\) 是 \(G\) 的正规 \(p\)-子群（不必是西罗子群）。

（a）证明 \(R\) 包含在 \(G\) 的每个西罗 \(p\)-子群中。

（b）若 \(S\) 是 \(G\) 的另一个正规 \(p\)-子群，证明 \(RS\) 也是 \(G\) 的正规 \(p\)-子群。

（c）子群 \(O_p(G)\) 定义为由 \(G\) 的所有正规 \(p\)-子群生成的群。证明 \(O_p(G)\) 是 \(G\) 的唯一的极大正规 \(p\)-子群，且 \(O_p(G)\) 等于 \(G\) 的所有西罗 \(p\)-子群的交。

（d）设 \(\bar{G} = G/O_p(G)\). 证明 \(O_p(\bar{G}) = 1\)（即 \(\bar{G}\) 没有非平凡正规 \(p\)-子群）。

\item 用西罗定理证明中的方法证明：若 \(n_p\) 不同余于 \(1 \pmod{p^2}\)，则存在 \(G\) 的不同西罗 \(p\)-子群 \(P\) 和 \(Q\) 使得 \(|P : P \cap Q| = |Q : P \cap Q| = p\).

\item 证明 \(GL_n(\mathbb{F}_p)\) 中严格上三角矩阵构成的子群（参见第 2.1 节习题 17）是这个有限群的西罗 \(p\)-子群。〔用第 1.4 节的阶的公式求 \(GL_n(\mathbb{F}_p)\) 的一个西罗 \(p\)-子群的阶。〕

\item 证明 \(GL_2(\mathbb{F}_p)\) 的西罗 \(p\)-子群的个数为 \(p+1\).〔给出两个不同的西罗 \(p\)-子群。〕

\item 证明 \(SL_2(\mathbb{F}_4) \cong A_5\)（关于 \(SL_2(\mathbb{F}_4)\) 的定义参见第 2.1 节习题 9）。

\item 证明 \(\mathbb{R}^3\) 中正二十面体的刚体运动群同构于 \(A_5\).〔回忆这个群的阶为 60：第 1.2 节习题 13。〕

\item 证明 \(\mathbb{R}^3\) 中正十二面体的刚体运动群同构于 \(A_5\).（与立方体和正四面体一样，正二十面体和正十二面体是互为对偶的立体。）〔回忆这个群的阶为 60：第 1.2 节习题 12。〕

\item 设 \(p\) 是整除有限群 \(G\) 的阶的最小素数。若 \(P \in \operatorname{Syl}_p(G)\) 且 \(P\) 循环，证明 \(N_G(P) = C_G(P)\).

\item 求 \(S_{2p}\) 的一个西罗 \(p\)-子群的生成元，其中 \(p\) 是奇素数。证明这是 \(p^2\) 阶阿贝尔群。

\item 求 \(S_{p^2}\) 的一个西罗 \(p\)-子群的生成元，其中 \(p\) 是素数。证明这是 \(p^{p+1}\) 阶非阿贝尔群。

\item 编写并执行一个计算机程序，它

（i）对每个小于 10000 的奇数 \(n\)，若 \(n\) 不是素数的幂且存在素因子 \(p\)，使得西罗定理的同余条件不能迫使 \(n_p\) 对 \(n\) 阶群都取 1，则输出该 \(n\)；

（ii）对（i）中每个 \(n\)，给出 \(n\) 的素数幂分解，并列出对整除 \(n\) 的每个素数 \(p\)，\(n_p\) 的所有允许取值（即那些不被西罗定理第 3 部分排除的取值）。

\item 对小于 1000 的所有偶数执行与前一习题相同的过程。解释列表的相对长度与所检验整数个数之间的关系。

\item 证明：若 \(|G| = 2^n m\)（其中 \(m\) 是奇数）且 \(G\) 有循环的西罗 2-子群，则 \(G\) 有 \(m\) 阶正规子群。〔使用归纳法和第 2 节习题 11 和 12。〕

\item 证明：若 \(U\) 和 \(W\) 是 \(G\) 的西罗 \(p\)-子群 \(P\) 的正规子集，则 \(U\) 在 \(G\) 中共轭于 \(W\) 当且仅当 \(U\) 在 \(N_G(P)\) 中共轭于 \(W\). 由此推出 \(P\) 的中心中的两个元素在 \(G\) 中共轭当且仅当它们在 \(N_G(P)\) 中共轭。（\(P\) 的子集 \(U\) 在 \(P\) 中正规是指 \(N_P(U) = P\).）

\item 设 \(P\) 是 \(G\) 的西罗 \(p\)-子群，\(M\) 是 \(G\) 的包含 \(N_G(P)\) 的任意子群。证明 \(|G : M| \equiv 1 \pmod{p}\).

\end{enumerate}

下面一串习题导向对具有下述性质的所有整数 \(n\) 的分类：每个 \(n\) 阶群都是循环群（例如 \(n = 15\) 就是这样一个整数）。这些论证是“每个奇数阶群都可解”这一证明的高度简化的原型，因为它们利用真子群的结构（交换性）及其在整个群中的嵌入（我们将看到不同的极大子群的交为单位元）通过计数论证得到矛盾。在奇数阶群可解的证明中，人们用归纳法化归到极小反例是单群的情形——但那里每个真子群都可解（不像我们的情形中是可交换的）。在这种情形下极大子群的结构与嵌入的分析要复杂得多，计数论证（大致说来）被特征标理论论证所取代（将在第 VI 部分讨论）。

\begin{enumerate}\setcounter{enumi}{51}

\item 假设 \(G\) 是每个真子群都阿贝尔的有限单群。若 \(M\) 和 \(N\) 是 \(G\) 的不同极大子群，证明 \(M \cap N = 1\).〔参见第 3 节习题 23。〕

\item 用前一习题证明：若 \(G\) 是每个真子群都阿贝尔的非阿贝尔群，则 \(G\) 不是单群。〔设 \(G\) 是该断言的极小反例，用第 3 节习题 24 表明 \(G\) 有多于一个极大子群的共轭类。用第 3 节习题 23 的方法计算落在 \(M\) 与 \(N\) 的所有共轭中的元素个数，其中 \(M, N\) 是 \(G\) 的不共轭极大子群；证明这给出多于 \(|G|\) 个元素。〕

\item 证明下面的分类：若 \(G\) 是 \(p_1p_2\cdots p_r\) 阶有限群，其中 \(p_i\) 是互不相同的素数，且对所有 \(i, j\) 都有 \(p_i \nmid p_j - 1\)，则 \(G\) 是循环群。〔由归纳法，\(G\) 的每个真子群都循环，故由前一习题 \(G\) 不是单群。若 \(N\) 是非平凡真正规子群，则 \(N\) 循环，且 \(G/N\) 作为 \(N\) 的自同构作用。用命题 \ref{pro:4.16} 证明 \(N \leq Z(G)\)，并用归纳法证明 \(G/Z(G)\) 循环，故由第 3.1 节习题 36，\(G\) 是阿贝尔群。〕

\item 证明前一习题的逆命题：若 \(n \geq 2\) 是使每个 \(n\) 阶群都循环的整数，则 \(n = p_1p_2\cdots p_r\) 是互不相同素数的乘积，且对所有 \(i, j\) 有 \(p_i \nmid p_j - 1\).〔若 \(n\) 不具此形式，用 \(p^2\) 阶和 \(pq\) 阶非循环群的直积构造 \(n\) 阶非循环群，其中 \(p \mid q - 1\).〕

\item 若 \(G\) 是每个真子群都阿贝尔的有限群，证明 \(G\) 可解。
\end{enumerate}

\section{\(A_n\) 的单性}

\(A_n\)（\(n \geq 5\)）单性的证明有很多种。最初等的一种是证明 \(A_n\) 由 3-循环生成，然后证明一个正规子群必含有一个 3-循环，从而必含有所有 3-循环，因此不能是真子群。我们给出一个计算量较小的途径。

注意 \(A_3\) 是阿贝尔单群，而 \(A_4\) 不是单群（\(n_3(A_4) = 1\)）。

\begin{theorem}\label{thm:4.24}
对所有 \(n \geq 5\)，\(A_n\) 是单群。
\end{theorem}

\begin{proof}
对 \(n\) 作归纳。结果已对 \(n = 5\) 建立，因此假设 \(n \geq 6\) 并令 \(G = A_n\). 假设存在 \(H \trianglelefteq G\)，满足 \(H \neq 1, G\).

对每个 \(i \in \{1, 2, \ldots, n\}\)，设 \(G_i\) 是 \(i\) 在 \(G\) 对 \(\{1, 2, \ldots, n\}\) 的自然作用下的稳定子。于是 \(G_i \leq G\) 且 \(G_i \cong A_{n-1}\). 由归纳假设，\(G_i\) 对 \(1 \leq i \leq n\) 都是单群。

首先假设存在 \(\tau \in H\) 使 \(\tau \neq 1\)，但对某个 \(i \in \{1, 2, \ldots, n\}\) 有 \(\tau(i) = i\). 由于 \(\tau \in H \cap G_i\) 且 \(H \cap G_i \trianglelefteq G_i\)，由 \(G_i\) 的单性必有 \(H \cap G_i = G_i\)，即
\[
G_i \leq H.
\]
由第 1 节习题 2，\(\sigma G_i \sigma^{-1} = G_{\sigma(i)}\)，故对所有 \(\sigma \in G\) 有 \(\sigma G_i \sigma^{-1} \leq \sigma H \sigma^{-1} = H\). 于是
\[
G_j \leq H, \quad \text{对所有 } j \in \{1, 2, \ldots, n\}.
\]
任意 \(\sigma \in A_n\) 可以写成偶数个（设为 \(2t\) 个）对换的乘积，即 \(\sigma = \lambda_1\lambda_2\cdots\lambda_t\)，其中每个 \(\lambda_k\) 是两个对换的乘积。由于 \(n > 4\)，每个 \(\lambda_k\) 落在某个 \(G_j\) 中，故
\[
G = \langle G_1, G_2, \ldots, G_n \rangle \leq H,
\]
这是矛盾。因此若 \(\tau \neq 1\) 是 \(H\) 的元素，则对所有 \(i \in \{1, 2, \ldots, n\}\) 有 \(\tau(i) \neq i\)，即 \(H\) 中没有非单位元固定 \(\{1, 2, \ldots, n\}\) 的任何元素。由此推出：若 \(\tau_1, \tau_2 \in H\) 且对某个 \(i\) 有 \(\tau_1(i) = \tau_2(i)\)，则
\[
\tau_1 = \tau_2, \tag{4.2}
\]
因为此时 \(\tau_2^{-1}\tau_1(i) = i\).

假设存在 \(\tau \in H\)，其循环分解中含有一个长度至少为 3 的循环，比如
\[
\tau = (a_1\ a_2\ a_3\ \ldots)(b_1\ b_2\ \ldots)\cdots .
\]
设 \(\sigma \in G\) 满足 \(\sigma(a_1) = a_1\)、\(\sigma(a_2) = a_2\) 但 \(\sigma(a_3) \neq a_3\)（注意这样的 \(\sigma\) 存在于 \(A_n\) 中，因为 \(n \geq 5\)）。由命题 \ref{pro:4.10}，
\[
\tau_1 = \sigma\tau\sigma^{-1} = (a_1\ a_2\ \sigma(a_3)\ \ldots)(\sigma(b_1)\ \sigma(b_2)\ \ldots)\cdots,
\]
故 \(\tau\) 和 \(\tau_1\) 是 \(H\) 的不同元素，且 \(\tau(a_1) = \tau_1(a_1) = a_2\)，与 (4.2) 矛盾。这证明了 \(H\) 的非单位元的循环分解中只能出现 2-循环。

设 \(\tau \in H\)，\(\tau \neq 1\)，故
\[
\tau = (a_1\ a_2)(a_3\ a_4)(a_5\ a_6)\cdots
\]
（注意这里用到了 \(n \geq 6\)）。设 \(\sigma = (a_1\ a_2)(a_3\ a_5) \in G\). 则
\[
\tau_1 = \sigma\tau\sigma^{-1} = (a_1\ a_2)(a_5\ a_4)(a_3\ a_6)\cdots,
\]
故 \(\tau\) 和 \(\tau_1\) 是 \(H\) 的不同元素，且 \(\tau(a_1) = \tau_1(a_1) = a_2\)，再次与 (4.2) 矛盾。这完成了 \(A_n\) 单性的证明。
\end{proof}

\subsection*{习题}
设 \(G\) 是群，\(\Omega\) 是无限集。

\begin{enumerate}
\item 证明对所有 \(n \geq 5\)，\(A_n\) 没有指数小于 \(n\) 的真子群。

\item 求对所有 \(n \geq 5\) 的 \(S_n\) 的所有正规子群。

\item 证明对所有 \(n \geq 5\)，\(A_n\) 是 \(S_n\) 中唯一的指数小于 \(n\) 的真子群。

\item 证明对每个 \(n \geq 3\)，\(A_n\) 由所有 3-循环组成的集合生成。

\item 证明：若 \(G\) 是一串单子群 \(G_1 \leq G_2 \leq \cdots \leq G\) 的并，即 \(G = \bigcup_{i=1}^{\infty} G_i\)，则 \(G\) 是单群。

\item 设 \(D\) 是 \(S_\Omega\) 中只移动 \(\Omega\) 的有限多个元素的置换组成的子群（在第 3 节习题 17 中描述），设 \(A\) 是所有这样的 \(\sigma \in D\) 的集合：\(\sigma\) 在它所移动的（有限）点集上起偶置换的作用。证明 \(A\) 是无限单群。〔证明 \(D\) 中每一对元素都落在 \(D\) 的一个有限单子群中。〕

\item 在上一习题的记号下，证明：若 \(H \trianglelefteq S_\Omega\) 且 \(H \neq 1\)，则 \(A \leq H\)，即 \(A\) 是 \(S_\Omega\) 的唯一的（非平凡）极小正规子群。

\item 在前两个习题的记号下，证明 \(|D| = |A| = |\Omega|\). 由此推出：若 \(S_\Omega \cong S_{\Omega'}\)，则 \(|\Omega| = |\Omega'|\).〔使用 \(D\) 由对换生成这一事实。你可以假设基数为 \(|\Omega|\) 的集合的可数并与有限直积的基数也是 \(|\Omega|\)。〕
\end{enumerate}
